\documentclass[12pt]{elsarticle}
\usepackage[utf8]{inputenc}
\usepackage[T1]{fontenc}
\usepackage[top=2.4cm,bottom=2.2cm,left=2.0cm,right=2.0cm]{geometry}
\usepackage[english]{babel}
\usepackage[table,dvipsnames]{xcolor}
\usepackage{graphicx,array,booktabs,longtable,multirow,tabularx}
\usepackage{amsmath,amssymb,amsthm,mathrsfs,bm}
\usepackage{float,placeins,caption,setspace,gensymb,wasysym,times}
\usepackage{needspace}
\usepackage{tikz,tikz-qtree}
\usetikzlibrary{arrows,arrows.meta,automata,positioning,calc,decorations.pathreplacing,decorations.pathmorphing,trees,shapes.misc}
\usepackage{etoolbox}
\usepackage[colorlinks=true,allcolors=black]{hyperref}
\urlstyle{same}
\AtBeginDocument{\hypersetup{allcolors=black}}
\definecolor{RevisionBlue}{RGB}{23,64,156}
\definecolor{RevisionRed}{RGB}{169,25,35}
\definecolor{nGray}{RGB}{214,214,214}
\definecolor{intnull2}{RGB}{214,214,214}
\newenvironment{revision}{\begingroup\color{RevisionBlue}\hypersetup{allcolors=RevisionBlue}}{\par\endgroup}
\DeclareRobustCommand{\rev}[1]{{\color{RevisionBlue}\hypersetup{allcolors=RevisionBlue}#1}}
\newenvironment{redrevision}{\begingroup\color{RevisionRed}\hypersetup{allcolors=RevisionRed}}{\par\endgroup}
\DeclareRobustCommand{\revred}[1]{{\color{RevisionRed}\hypersetup{allcolors=RevisionRed}#1}}
\pdfstringdefDisableCommands{\def\rev#1{#1}\def\revred#1{#1}}
\AtBeginEnvironment{thebibliography}{%
  \let\savedbibitem\bibitem
  \renewcommand{\bibitem}[2][]{%
    \color{black}\hypersetup{allcolors=black}%
    \ifboolexpr{test {\ifstrequal{#2}{zhang2025myoelectric}} or test {\ifstrequal{#2}{zhang2024fall}} or test {\ifstrequal{#2}{chen2022computer}} or test {\ifstrequal{#2}{sethi2022comprehensive}} or test {\ifstrequal{#2}{scataglini2024accuracy}} or test {\ifstrequal{#2}{boswell2021neural}}}{\color{RevisionBlue}\hypersetup{allcolors=RevisionBlue}}{}%
    \ifboolexpr{test {\ifstrequal{#2}{wade2023examination}} or test {\ifstrequal{#2}{uhlrich2023opencap}} or test {\ifstrequal{#2}{augustine2025clothing}} or test {\ifstrequal{#2}{viswakumar2022omgait}} or test {\ifstrequal{#2}{lee2023lowlight}} or test {\ifstrequal{#2}{wu2012illumination}}}{\color{RevisionRed}\hypersetup{allcolors=RevisionRed}}{}%
    \savedbibitem[#1]{#2}}}
\newcommand{\etal}{\textit{et al}}
\newcommand{\etc}{\textit{etc.}}
\newcommand{\citeColored}[2]{\cite{#2}}
\setstretch{1}
\setlength{\emergencystretch}{2em}
\captionsetup[table]{position=top,skip=7pt}
\setlength{\LTpre}{7pt}
\setlength{\LTpost}{7pt}
\setlength{\LTcapwidth}{\textwidth}
\journal{Neural Networks}

\renewcommand{\thesection}{S\arabic{section}}
\renewcommand{\thetable}{S\arabic{table}}
\begin{document}
\begin{center}
{\Large\bfseries Supplementary Material}\\[8pt]
{\large Vision-Based Gait Analysis for Identifying Lower-Limb Joint Mobility Disorders: A Comprehensive Survey}
\end{center}
\begin{revision}
This supplement contains the detailed study descriptions and extended comparison tables supporting the shorter synthesis in the main manuscript. Material relocated without a substantive change is shown in black. New explanations and revised wording are shown in blue. The study summaries retain the measurements reported in the original manuscript; methodological transfer is not treated as clinical validation. References are numbered separately within this supplement.
\end{revision}
\tableofcontents
\clearpage
\section{Hip: study details}
\label{supp:S1}
{Kiss et al. \cite{kiss2010effect} examined how gait speed and severity influence gait variability in hip OA. Variability was quantified as the standard deviation of spatio-temporal parameters and the mean standard deviation of angular parameters. The cohort comprised 20 patients with moderate hip OA, 20 with severe hip OA and 20 healthy controls, all assessed with the HHS and recorded with the Zebris CMS-HS at a self-selected speed and at treadmill speeds of 0.8, 1.0 and 1.2~m/s. Asymmetry was defined as the difference between affected and unaffected sides in patients and between dominant and non-dominant sides in controls. HHS was significantly better in moderate than in severe OA, and self-selected speed differed significantly across groups. Controls showed the lowest spatio-temporal variability and severe OA the highest. Severity mainly affected the variability of hip and knee motion and pelvic movement, indicating less stable and more irregular gait with advanced disease.}

{ Shakoor et al. \cite{shakoor2003asymmetric} studied why the contralateral knee is often at greater risk of degeneration than the ipsilateral knee in unilateral hip OA. Fifty patients awaiting unilateral total hip replacement were assessed preoperatively with marker-based multi-camera gait analysis and force platforms at self-selected speed, and 22 were reassessed one to two years after surgery. Using inverse dynamics and a predictive knee loading model, the contralateral knee showed significantly greater external adduction, extension and internal rotational moments and predicted medial compartment load than the ipsilateral knee before surgery. After surgery, hip pain was substantially reduced, yet the contralateral knee still showed higher adduction moment, medial compartment load and extension moment. Hip pathology therefore alters loading across the lower limb and may contribute to secondary joint deterioration.}

{ Ismailidis et al. \cite{ismailidis2021kinematic} investigated the effect of walking speed on spatio-temporal and kinematic characteristics in 24 patients with severe unilateral hip OA and 48 healthy asymptomatic controls. Symptoms and function were assessed with the HOOS questionnaire and gait was measured with the RehaGait wearable sensor system. Patients walked at a self-selected speed, and controls were analysed at a normal speed and at speed-matched conditions. Patients had significantly lower HOOS scores in all subdomains, and walked more slowly with lower cadence than controls. After speed matching, some spatio-temporal differences became less pronounced, but clear kinematic differences remained: reduced hip flexion in stance and swing, reduced hip extension in late stance, increased knee flexion in late stance, and greater ankle dorsiflexion in early swing and late stance. These deficits cannot be explained solely by slower walking and reflect disease-specific movement adaptations.}

{ Laroche et al. \cite{laroche2014classification} explored whether patients with hip OA could be distinguished from healthy controls using three-dimensional gait trajectories and machine learning. Twenty symptomatic patients with stage 2 or 3 hip OA and 20 controls walked ten trials along a 4-m walkway, recorded by an eight-camera marker-based system. Angular trajectories of the foot, shank, thigh, trunk and hip were extracted in three planes. A support vector machine trained on full gait-cycle time series, with five-fold cross-validation, reached 88\% accuracy using all kinematic trajectories. Receiver operating characteristic-based cycle classification across multiple gait cycles increased accuracy to as high as 97\%. The sagittal thigh angle alone achieved 85.8\%. The maximum amplitudes of these trajectories correlated strongly with WOMAC function scores, linking classification performance to clinical severity.}

{ Meyer et al. \cite{meyer2015biomechanical} also addressed hip OA classification, aiming to identify a compact set of clinically meaningful variables. Twenty patients with hip OA and 17 controls were assessed with a 15-camera Vicon system, 22 reflective markers and two force platforms. Twenty-seven kinematic and kinetic waveforms were retained, reduced by principal component analysis (PCA) and classified with linear discriminant analysis (LDA). Patients walked more slowly, with shorter stride length and longer double-support time. Of the 44 principal components retained, the largest contributions came from the hip, knee, ankle and trunk/pelvis segments. Models using up to ten components gave classification errors of 15\% to 19\%, corresponding to approximately 81\% to 85\% accuracy. Hip OA can therefore be identified from a small number of interpretable gait features.}

{ Reininga et al. \cite{reininga2012subjects} focused on frontal-plane trunk compensation in hip OA, such as the Duchenne limp, in which patients modify trunk movement to offload the affected hip. Sixty patients with hip OA, of whom 10 had a clear Duchenne limp, and 30 controls walked along a 25-m corridor at self-selected speed. Two triaxial body-fixed units incorporating gyroscopes, accelerometers and magnetometers measured pelvic and thoracic motion. Controls showed increasing pelvic range of motion (ROM) and decreasing thoracic ROM with speed, giving a thoracic-to-pelvic ROM ratio below one. Patients with Duchenne limp had thoracic ROM greater than pelvic ROM at all speeds, giving a ratio above one. Patients without the limp had reduced pelvic ROM and a ratio close to one. Both patient groups walked more slowly than controls, showing that trunk and pelvic motion reveal compensatory strategies in hip OA.}

{ Eitzen et al. \cite{eitzen2012sagittal} examined gait in 48 patients with mild-to-moderate hip OA who were not eligible for hip replacement, and 22 controls. \rev{The HHS described clinical status; radiographic severity was considered separately.} Gait was captured with an eight-camera marker-based system and force platforms at self-selected speed. The OA group had significantly reduced gait speed, reduced hip ROM in most directions, lower knee extensor strength, reduced hip and knee excursions, and lower hip and knee extension in late stance. In the subgroup with more severe radiographic narrowing, hip excursion was lower, midstance knee extension was reduced and peak hip flexion moment was lower. Patients with relatively mild symptoms can therefore show clear biomechanical deficits, suggesting gait analysis may be useful before surgery becomes necessary.}

{ Mont et al. \cite{mont2007gait} compared gait across patients with resurfacing total hip arthroplasty (RTHA), standard total hip arthroplasty (STHA), hip OA and healthy controls. The study included 15 participants in each arthroplasty group, 10 patients with hip OA and 30 matched controls. Gait was analysed with eight cameras, two force platforms and 26 reflective markers. RTHA patients walked faster than STHA and OA patients, at speeds similar to controls. Hip abductor and extensor moments in the RTHA group were comparable to controls but reduced in the OA and STHA groups. Knee flexion and ankle dorsiflexion in the arthroplasty groups were closer to normal than in the OA group. Well-functioning resurfacing arthroplasty \rev{showed favourable gait findings in this cohort; the group comparison does not establish superiority for all patients.}}

The existing literature shows that hip OA affects gait in a { complex,} multi-level manner, influencing the { hip, knee, ankle, pelvis and trunk. The} most consistent findings include reduced walking speed, altered hip and knee kinematics, increased gait variability, abnormal { asymmetry} and compensatory upper-body motion \cite{kiss2010effect,ismailidis2021kinematic,reininga2012subjects,eitzen2012sagittal,meyer2015biomechanical}. { Gait analysis supports} several clinical goals, including severity assessment, identification of compensatory mechanisms, evaluation of surgical { outcomes and data-driven patient} classification \cite{laroche2014classification,meyer2015biomechanical,mont2007gait}. This makes gait analysis relevant { to both clinical biomechanics and} the development of computer-vision-based systems for automated hip mobility assessment.

{\scriptsize
\setlength{\tabcolsep}{2pt}
\renewcommand{\arraystretch}{1.03}
\begin{longtable}{p{1.9cm} p{1.6cm} p{1.9cm} p{1.9cm} p{1.7cm} p{1.9cm} p{2.0cm} p{3.0cm}}
\caption{Critical comparison of existing literature on hip-related mobility disorders and gait analysis.}
\label{hip_critical_extended_summary} \\
\hline
Reference & Cohort & Focus & Data / Setup & Analysis & Contribution & Strength & Limitation \\
\hline
\endfirsthead

\hline
Reference & Cohort & Focus & Data / Setup & Analysis & Contribution & Strength & Limitation \\
\hline
\endhead

Kiss \cite{kiss2010effect}
& 20 moderate OA, 20 severe OA, 20 controls
& Speed and severity effects
& Zebris CMS-HS; treadmill; HHS
& ANOVA; SPSS
& Linked OA severity and gait speed with variability
& Clear severity-based comparison with clinical scoring
& Small cohort; treadmill gait may not reflect natural walking; mainly descriptive \\
\hline

Shakoor et al. \cite{shakoor2003asymmetric}
& 50 unilateral hip OA; 22 post-THR follow-up
& Contralateral knee loading
& Marker-based gait; force mat; HHS
& Inverse dynamics; load modelling
& Showed persistent contralateral knee overloading
& Strong translational value across joints
& Reduced follow-up sample; not designed for hip classification \\
\hline

Ismailidis et al. \cite{ismailidis2021kinematic}
& 24 severe unilateral OA, 48 controls
& Speed-matched deficits
& RehaGait; HOOS
& SPSS; SPM
& Identified disease-specific kinematics beyond slower speed
& Strong control for speed confounding
& Severe unilateral OA only; wearable detail lower than optical systems \\
\hline

Laroche et al. \cite{laroche2014classification}
& 20 stage 2--3 OA, 20 controls
& Hip OA classification
& 8-camera gait; WOMAC
& SVM; ROC; ANOVA
& High discrimination using gait-cycle trajectories
& Strong proof-of-concept for gait-based diagnosis
& Small sample; no external validation; laboratory-dependent pipeline \\
\hline

Meyer et al. \cite{meyer2015biomechanical}
& 20 OA, 17 controls
& Low-dimensional classification
& 15-camera Vicon; force plates
& PCA; LDA
& Reduced gait space to compact discriminative variables
& Good balance of reduction and interpretation
& Small cohort; specialised setup; moderate classification error \\
\hline

Reininga et al. \cite{reininga2012subjects}
& 60 hip OA, 30 controls
& Trunk compensation
& 2 triaxial sensors
& Xsens; PASW
& Quantified Duchenne limp and frontal-plane compensation
& Clinically meaningful and more practical than full lab systems
& Limited joint-level detail; not a classification study \\
\hline

Eitzen et al. \cite{eitzen2012sagittal}
& 48 mild-moderate OA, 22 controls
& Early-stage deficits
& 8-camera gait; force plates; HHS; WOMAC
& Visual3D; SPSS
& Showed measurable deficits before surgery eligibility
& Valuable for early disease assessment
& Subgroup differences were modest; limited predictive use \\
\hline

Mont et al. \cite{mont2007gait}
& 15 RTHA, 15 STHA, 10 OA, 30 controls
& Surgical comparison
& 8 cameras; 2 force mats; 26 markers
& OrthoTrak
& Compared gait restoration after different arthroplasties
& Strong clinical comparison across groups
& Small surgical cohorts; not a predictive framework \\
\hline

Aminian et al. \cite{aminian2004evaluation}
& Hip OA and post-THR
& Ambulatory evaluation
& Mobile gait analysis system
& Comparative evaluation
& Demonstrated feasibility beyond gait labs
& Supports scalable monitoring
& Focused on system validation more than disease modelling \\
\hline

Hodt-Billington et al. \cite{hodt2012criteria}
& Hip OA patients
& Gait asymmetry
& Clinical gait asymmetry assessment
& Statistical comparison
& Showed asymmetry as a useful functional marker
& Clinically intuitive
& Narrow biomechanical scope; no predictive model \\
\hline

Ismailidis et al. \cite{ismailidis2020measuring}
& Severe hip OA
& Wearable kinematics
& Wearable sensors
& Sensor-based kinematic analysis
& Supported practical monitoring of hip OA gait
& Scalable and repeatable
& Lower anatomical precision than optical motion capture \\
\hline

\end{longtable}
}

Table~\ref{hip_critical_extended_summary} provides a critical comparison of the literature studies on hip-related mobility disorders and gait analysis. Taken together, these studies span a broad range of clinical scenarios, including hip osteoarthritis, post-arthroplasty assessment, compensatory gait behaviour, and severity-based subgroup analysis. They also reflect a wide variety of methodological choices, from camera-based motion capture and force platforms to wearable and ambulatory systems, combined with analytical approaches such as statistical testing, inverse dynamics, principal component analysis (PCA)/linear discriminant analysis (LDA), support vector machine (SVM), receiver operating characteristic (ROC) analysis, and statistical parametric mapping (SPM). The comparison shows that hip-related disorders consistently affect gait through changes in walking speed, asymmetry, compensatory movement, joint excursions, and cross-joint loading patterns. It also suggests that gait-based classification and discrimination of hip pathology are feasible, particularly when clinically meaningful biomechanical features are combined with structured analytical models. An important strength of the existing literature is that it captures both mechanistic and diagnostic perspectives, offering insight not only into how hip pathology alters movement, but also into how these alterations may support assessment, stratification, and outcome evaluation. Also, Table~\ref{hip_critical_extended_summary} makes clear that the field still faces several recurring limitations. Much of the evidence is based on relatively small cohorts, laboratory-constrained protocols, inconsistent feature definitions, modest external validation, and limited explainability or scalability. These issues reduce clinical transferability and make it difficult to compare findings across studies in a fully consistent way. For this reason, this type of comparison helps to position the need for more robust, interpretable, and scalable frameworks for hip mobility assessment, particularly those that can retain biomechanical relevance while being more practical for wider clinical and real-world use.



\section{Knee: study details}
\label{supp:S2}
Simonsen et al. \cite{simonsen2012gait} investigated adults with { generalised} joint hypermobility (GJH), a condition { of} increased joint range of motion { that} may alter loading and { raise} the long-term risk of { OA. Seventeen} participants with GJH and 17 without { GJH were included. GJH was defined as} a Beighton score { above} 4 together with at least one hypermobile knee. Gait was analysed on a { 10-m walkway with} five marker-based cameras and force platforms, { and three-dimensional coordinates were reconstructed by} direct linear transformation. { Knee pain in} the preceding week was substantially more common in the GJH group. { Participants with GJH showed} a smaller frontal hip angle and greater { sagittal} knee flexion at heel strike and during early stance. { They also had} higher joint moments in the frontal and sagittal planes, suggesting increased knee loading and a possible { future risk of OA.}

Bejek et al. \cite{bejek2006influence} examined the effect of walking speed on gait parameters in knee and hip { OA,} compared with healthy older adults. { Speed} is a major determinant of many gait variables and must be considered when comparing patient and control groups. The study included patients with severe knee OA, patients with severe hip { OA} and healthy elderly { controls, who walked on a treadmill at different speeds. The} highest speed at which the OA groups could walk without pain and with acceptable bilateral { coordination was 2.00~km/h. In controls,} 17 of 22 gait parameters changed with speed, { and in knee OA, 15 parameters changed.} Gait asymmetry was greater in both OA groups than in controls. { Changes} in gait may { therefore} reflect both pathology and walking speed, and speed effects should be accounted for when interpreting OA-related adaptations.

Kiss \cite{kiss2011effect} { studied} how knee OA { severity} influences gait variability, { which} may indicate reduced stability and a higher risk of falling. The study compared 45 patients with moderate OA (Kellgren--Lawrence grade 3), 45 with severe OA (grade { 4) and healthy controls,} assessed on a treadmill at 1.0, { 1.1} and 1.2~m/s. Variability was quantified { by} the coefficient of variation (CV) for spatio-temporal and angular parameters. { OA severity clearly influenced variability.} Angular features of the knee and hip, { and} pelvic tilt and rotation, were { most affected, while} several spatio-temporal variables were less sensitive. { As severity increased, spatio-temporal variability rose but angular variability fell. Females showed} greater spatio-temporal variability than males. { Disease severity therefore} shapes the consistency and stability of { movement, not only mean gait behaviour.}

Elbaz et al. \cite{elbaz2014novel} proposed an objective gait-based system for grading knee OA severity using only spatio-temporal parameters. { The} study included 2,911 patients with knee OA, all of whom underwent { computerised gait analysis, Kellgren--Lawrence radiographic grading} and WOMAC and { SF-36 questionnaires. Because} step length and body { size differ by sex, separate models were built} for men and women. { The classification relied mainly} on stride length and cadence and achieved accuracies of 90.8\% for { men} and 89.5\% for { women. The approach aims for} a functionally grounded { grading} based on real walking { performance, rather than imaging alone.}

Bacon et al. \cite{bacon2022relation} examined the relationship between gait and mild unilateral knee pain using MOST { data.} After excluding participants with bilateral pain, 2,066 individuals who had completed a 20-m gait { test were analysed.} Machine learning first { identified} gait parameters linked to unilateral pain. { In} mild unilateral pain, lower step regularity and lower sample entropy were associated with { a} greater probability of pain, whereas gait speed was not the key factor. In moderate-to-severe unilateral pain, reduced gait speed was associated with { a higher probability of pain. Early pain-related gait changes may therefore appear as} reduced regularity before slowing becomes apparent.

{ Qingguang et al. \cite{qingguang2015gait} evaluated Chinese massage therapy in 20 women with knee OA. Treatment was given three times per week for two weeks, with sessions of 25--30 minutes. Marker-based three-dimensional gait analysis and WOMAC were recorded before and after treatment. WOMAC scores improved significantly, indicating less pain and stiffness and better function. Walking speed, step width and support time also improved, and sagittal range of motion increased at the hip, knee and ankle, although initial contact angles changed little. The sample was small, but the study suggests gait analysis can quantify short-term functional gains after conservative treatment.}

Ro et al. \cite{ro2019effects} examined whether knee OA affects hip and ankle { mechanics as well as the knee. Eighty-nine} patients with end-stage knee OA and 42 { controls} underwent gait analysis and radiographic evaluation. { Hip} and ankle range of motion { were} significantly reduced in the OA { group} regardless of gait speed. { Sagittal} and coronal motion arcs were { smaller} and correlated with knee range of motion. { Sagittal hip, ankle} and knee moments were { reduced,} while coronal ankle and knee moments were increased. { Knee OA therefore} has system-wide biomechanical consequences, and secondary changes at the hip and { ankle deserve attention.}

Butler et al. \cite{butler2009gait} { analysed} gait after anterior cruciate ligament (ACL) { reconstruction, since ACL injury increases} the risk of early-onset knee { OA.} Seventeen individuals with prior reconstruction were compared with healthy controls during controlled-speed { walking,} using force platforms and eight marker-based cameras. The ACL group { had} a 21\% higher peak internal knee abduction moment, a feature linked { to} medial knee loading and OA progression. { Hip moments did not differ between groups. Altered} knee loading may { therefore persist after reconstruction} and contribute to later { degeneration.}

Harding et al. \cite{harding2012body} investigated the interaction between BMI and gait in moderate knee OA. { The} study included 104 asymptomatic participants and 140 with moderate OA, both groups stratified into healthy-weight, { overweight} and obese subgroups. Marker-based analysis and force platform data were { combined} with principal component analysis { (PCA)} to evaluate knee moment patterns. { Participants} with OA showed reduced knee flexion and extension angles, larger knee moments and altered rotational mechanics. { Within} the OA group, higher BMI was associated with further reductions in knee flexion range and altered moment patterns during stance. { Obesity therefore modifies how} patients with OA move, potentially amplifying { impairment rather than simply adding mechanical load.}

Da Silva et al. \cite{da2012biomechanical} { quantified} biomechanical change in medial knee { OA, comparing patients with asymptomatic controls using} a six-camera motion capture system and force platform { data. The} Knee Society Score { was used} to capture clinical function. The OA group { had} significantly lower cadence, step { length} and walking velocity, while { stance-phase percentage changed less.} Peak knee flexion and extension angles were slightly reduced, but moments differed more clearly: increased peak { moments and} external rotation moments, and reduced internal rotation moments. The authors { attributed these changes to} pain-related adaptation and altered loading strategies.

Kaufman et al. \cite{kaufman2001gait} { analysed} level walking, stair { ascent and stair descent in} 139 adults with knee OA and 20 healthy { controls, using a six-camera system. The} OA group walked more slowly and had a significantly reduced peak knee extension moment. Knee flexion was similar between groups during level walking, { with} task-specific differences during stair negotiation. { Women} showed greater peak knee flexion, which may partly reflect anthropometric differences, and higher BMI was associated with lower knee extension moments. { Patients} adapt their gait to reduce pain and joint stress, but { these} adaptations may reduce functional efficiency.

Weidow et al. \cite{weidow2006hip} examined whether gait mechanics differ between medial and lateral knee OA. { Women} with medial OA, women with lateral { OA} and healthy female { controls were assessed with} six infrared cameras and force platforms. Both OA groups walked more slowly and with shorter stride length than { controls.} The medial OA group { had a} reduced peak knee extension angle and a markedly higher peak knee abduction { moment. The} lateral OA group { had a} reduced knee flexion angle and a much lower knee abduction moment relative to controls. { Femoral} and tibial rotation also { differed,} with increased external femoral rotation { in lateral OA} and increased tibial { rotation in medial OA. Medial} and lateral OA should not be treated as mechanically { identical.}

Favre and Jolles \cite{favre2016gait} reviewed { recurring gait hallmarks of knee OA:} increased knee adduction moment (KAM), reduced knee flexion moment { (KFM)} and reduced knee flexion angle (KFA). KAM was consistently higher in OA patients than in controls and tended to be greater in more severe disease. KFM was often reduced, suggesting a pain-avoidance strategy that limits sagittal-plane loading. Reduced KFA reflected decreased knee range of motion. { The review distils these} findings into clinically meaningful biomechanical indicators.

Phinyomark et al. \cite{phinyomark2016gender} examined sex-related differences in gait among patients with knee OA. { The} study included 45 men and 55 women with { OA and} 43 healthy individuals, recorded on a treadmill { with} an eight-camera three-dimensional system. { Discrete} variable analysis and { PCA} were applied to the kinematic data. { Women with OA had} greater knee abduction angle and peak hip adduction than men with OA, while overall range of motion was similarly reduced in both groups relative to controls. { Of} the 112 discrete variables examined, significant male--female differences were found mainly in frontal-plane hip and knee motion rather than in sagittal-plane measures. { Sex therefore} influences specific movement patterns in OA and may { matter for personalised interventions and} predictive models.

B{\k{a}}czkowicz et al. \cite{bkaczkowicz2018gait} assessed gait recovery after total knee arthroplasty { (TKA),} combining spatio-temporal analysis with { WOMAC patient-reported outcomes. Twenty-one patients were assessed} before surgery, five days { and} 12 weeks after surgery, and after rehabilitation, with 60 healthy participants as controls. Gait speed, step { length} and single-support time were measured { with a gait mat. WOMAC scores fell} by 40\% before rehabilitation and by 55\% after { rehabilitation,} relative to the preoperative baseline. Gait speed initially declined after { surgery and then improved towards} control levels. Step length asymmetry { decreased} and support time { became more balanced} between the operated and non-operated { legs.} The negative association between gait speed and WOMAC suggests that objective gait recovery parallels patients' { perceived} improvement.

Astephen et al. \cite{astephen2008biomechanical} { analysed} how knee OA severity affects gait across the hip, { knee and ankle,} using the Optotrak system and force { platforms in healthy, moderate and severe OA groups. Thirty-four gait variables were derived} from 12 underlying kinetic and kinematic measures. { Walking} speed decreased progressively from healthy controls to moderate and then severe OA, { and} stride length was particularly reduced in severe disease. Some changes were common to all OA patients, { including} reduced early-stance flexion moment and increased mid-stance adduction { moment. Others} were specific to severe OA, including reduced knee extension moment, lower late-stance internal rotation { moment} and reduced sagittal-plane motion at the hip, { knee} and ankle. Progressive changes were also { seen} in knee flexion angle and { stance extension moment. Severity is therefore} reflected across the entire lower-limb kinetic chain.

Finally, McClelland et al. \cite{mcclelland2007gait} reviewed gait studies { after} total knee replacement (TKR) { across 11 studies. Most} patient groups walked more slowly than controls, and reduced knee range of motion { was} common even after surgery. Lower knee flexion during swing and reduced flexion during loading response were also { frequent. A} biphasic knee moment pattern was seen in around 80\% of healthy controls { but} in only 20\%--36\% of TKR patients. { Although} TKR can reduce pain and improve function, recovery of normal gait mechanics is often incomplete.

In summary, the reviewed literature presents a consistent picture: knee mobility disorders alter gait across multiple levels, from basic spatio-temporal measures to joint moments and inter-joint { coordination.} Gait analysis has been used to study pain, severity, phenotype differences, obesity-related effects, postoperative { recovery} and early risk states such as ACL reconstruction and hypermobility \cite{simonsen2012gait,butler2009gait,harding2012body,astephen2008biomechanical}. { Several} studies show that gait-based features can support classification and severity assessment in a clinically meaningful and potentially objective { way} \cite{elbaz2014novel,bacon2022relation,weidow2006hip}. This makes gait analysis highly relevant { to clinical biomechanics and to} the development of { intelligent,} computer-vision-based systems for automated knee mobility assessment.




{\scriptsize
\setlength{\tabcolsep}{2pt}
\renewcommand{\arraystretch}{1.03}
\begin{longtable}{p{1.9cm} p{1.6cm} p{1.9cm} p{1.9cm} p{1.7cm} p{1.9cm} p{2.0cm} p{3.0cm}}
\caption{Critical comparison of existing literature on knee-related mobility disorders and gait analysis.}
\label{knee_critical_extended_summary} \\
\hline
Reference & Cohort & Focus & Data / Setup & Analysis & Contribution & Strength & Limitation \\
\hline
\endfirsthead

\hline
Reference & Cohort & Focus & Data / Setup & Analysis & Contribution & Strength & Limitation \\
\hline
\endhead

Simonsen et al. \cite{simonsen2012gait}
& 17 GJH, 17 non-GJH
& Early biomechanical risk
& 5-camera marker-based gait; force platforms
& 3D reconstruction; MATLAB-based analysis
& Showed higher knee loading and altered sagittal/frontal mechanics in GJH
& Valuable for identifying pre-osteoarthritic biomechanical risk
& Small cross-sectional cohort; does not show whether altered loading leads to later OA \\
\hline

Bejek et al. \cite{bejek2006influence}
& Knee OA, hip OA, healthy controls
& Effect of gait speed
& Treadmill gait analysis
& Comparative gait parameter analysis
& Demonstrated that walking speed strongly alters many gait variables
& Important methodological reminder that speed is a major confounder
& Focuses more on speed sensitivity than disease-specific mechanisms; treadmill gait may differ from natural walking \\
\hline

Kiss \cite{kiss2011effect}
& 45 moderate OA, 45 severe OA, controls
& Severity and variability
& Treadmill gait at 3 speeds
& Coefficient of variation; SPSS
& Linked OA severity with altered gait variability patterns
& Moves beyond mean values and captures gait stability
& Laboratory treadmill design may constrain natural gait; interpretation of opposite variability trends needs further validation \\
\hline

Elbaz et al. \cite{elbaz2014novel}
& 2911 knee OA patients
& Objective OA grading
& Computerised gait analysis; WOMAC; SF-36; radiographic grading
& Rule-based classification
& Proposed a simple gait-based system for grading knee OA
& Very large cohort and clinically practical variables
& Relies mainly on stride length and cadence; richer kinematic and kinetic information is not exploited \\
\hline

Bacon et al. \cite{bacon2022relation}
& 2066 MOST participants with unilateral or no pain
& Pain-related gait signatures
& Cohort gait data from MOST
& ML feature selection; covariate-adjusted analysis
& Showed that mild unilateral pain is linked more to gait regularity than to speed
& Large population-based dataset capturing subtle early changes
& Focuses on pain association rather than full biomechanical phenotype; observational design limits causal claims \\
\hline

Qingguang et al. \cite{qingguang2015gait}
& 20 women with knee OA
& Treatment response
& 3D marker-based gait; WOMAC; massage intervention
& Pre-post statistical analysis
& Reported improved gait and symptoms after Chinese massage therapy
& Useful example of gait analysis for short-term treatment evaluation
& Small sample, no strong control group, and short follow-up reduce confidence in long-term impact \\
\hline

Ro et al. \cite{ro2019effects}
& 89 end-stage knee OA, 42 controls
& Whole-limb effects
& Marker-based gait analysis; radiographic evaluation
& Kinematic and kinetic comparison
& Showed that knee OA alters hip and ankle mechanics as well
& Strong evidence that knee OA affects the full lower-limb chain
& Focus on end-stage OA limits generalisability to earlier disease \\
\hline

Butler et al. \cite{butler2009gait}
& 17 ACL-reconstructed participants and controls
& Pre-OA risk after ACL reconstruction
& 8-camera marker-based gait; force platforms
& Group comparison of OA-related gait metrics
& Identified higher knee abduction moment after ACL reconstruction
& Extends OA-related gait analysis to an at-risk population
& Small sample and narrow variable set limit broader biomechanical interpretation \\
\hline

Harding et al. \cite{harding2012body}
& 140 moderate OA, 104 controls stratified by BMI
& BMI and OA interaction
& Marker-based gait; force platforms
& PCA-based moment analysis
& Showed that higher BMI amplifies OA-related gait abnormalities
& Important for linking obesity with altered knee mechanics
& Mainly moderate OA cases; effects across the full severity spectrum remain unclear \\
\hline

da Silva et al. \cite{da2012biomechanical}
& Medial knee OA and controls
& Medial OA biomechanics
& 6-camera gait analysis; force platforms; KSS
& SPSS-based comparison
& Characterised altered cadence, step length, velocity, and rotational moments
& Useful for medial OA-specific description
& Modest analytical depth and limited sample detail restrict broader inference \\
\hline

Kaufman et al. \cite{kaufman2001gait}
& 139 knee OA, 20 controls
& Functional gait adaptation
& 6-camera gait analysis during level walking and stairs
& Task-based biomechanical comparison
& Showed slower gait and reduced knee extension moment in OA
& Functionally relevant because it includes stair tasks
& Older protocol and smaller control group make comparison with newer work less direct \\
\hline

Weidow et al. \cite{weidow2006hip}
& 15 medial OA, 15 lateral OA, 15 controls
& Medial vs lateral OA
& 6-camera infrared gait; force platforms
& Detailed rotational and moment analysis
& Showed distinct biomechanical patterns in medial and lateral OA
& Important for demonstrating phenotype-specific gait differences
& Small female-only cohort limits broader generalisability \\
\hline

Favre and Jolles \cite{favre2016gait}
& Review of knee OA gait biomechanics
& Biomechanical hallmarks
& Review of kinematic and kinetic evidence
& Narrative synthesis
& Identified recurrent markers such as higher KAM and lower KFM
& Provides a clear biomechanical framework for interpretation
& As a review, conclusions depend on heterogeneous primary studies and do not resolve protocol inconsistency \\
\hline

Phinyomark et al. \cite{phinyomark2016gender}
& 45 males, 55 females with OA, plus controls
& Sex differences
& 8-camera treadmill gait analysis
& Discrete-variable and PCA analysis
& Showed sex-related differences mainly in frontal-plane hip and knee motion
& Highlights the value of sex-aware gait analysis
& Severe OA excluded; treadmill setup may not represent overground gait \\
\hline

B{\k{a}}czkowicz et al. \cite{bkaczkowicz2018gait}
& 21 TKA patients, 60 controls
& Recovery after arthroplasty
& Gait mat; WOMAC; repeated follow-up
& Longitudinal ANOVA-based comparison
& Tracked improvement in gait symmetry, speed, and self-reported function
& Useful longitudinal picture of post-surgical recovery
& Small surgical cohort and limited biomechanical variables compared with motion-capture studies \\
\hline

Astephen et al. \cite{astephen2008biomechanical}
& 60 controls, 60 moderate OA, 61 severe OA
& Severity-dependent biomechanics
& Optotrak motion capture; force platforms; WOMAC
& One-way ANOVA across 34 gait variables
& Showed progressive changes across hip, knee, and ankle with OA severity
& One of the strongest whole-limb severity studies
& Large variable set improves coverage but reduces immediate clinical interpretability \\
\hline

McClelland et al. \cite{mcclelland2007gait}
& Systematic review of post-TKR gait studies
& Outcomes after knee replacement
& Review of gait-based kinetic and kinematic studies
& Methodological synthesis
& Showed that gait often remains abnormal after TKR despite clinical improvement
& Important for distinguishing symptom relief from biomechanical recovery
& Conclusions are affected by heterogeneity in protocols, follow-up times, and reported variables \\
\hline

\end{longtable}
}

Table~\ref{knee_critical_extended_summary} offers a critical overview of representative studies on knee-related mobility disorders and gait analysis. The reviewed literature spans a diverse set of cohorts, including patients with osteoarthritis, individuals with generalized joint hypermobility, ACL-reconstructed populations, post-arthroplasty cases, and groups defined by disease severity. Collectively, these studies have used gait analysis to investigate a wide range of clinically meaningful questions, including diagnosis, severity stratification, biomechanical risk, treatment response, recovery patterns, and phenotype-specific differences. A message emerging from this comparison is that knee-related disorders affect gait in a systematic way, particularly through changes in walking speed, symmetry, joint loading, and coordination across multiple joints. The literature also shows that gait-based measures are useful not only for describing biomechanical dysfunction, but also for supporting grading, subgroup discrimination, and classification. One of the key strengths is that it combines biomechanical insight with clinical application, thereby improving understanding of both the mechanisms of altered movement and their practical value for assessment and monitoring. However, the comparison also reveals a number of recurring limitations. Much of the current evidence remains rooted in laboratory-based protocols, relatively small study populations, limited external validation, and gait features that are not always defined or reported consistently across studies. These issues make it harder to compare findings directly and reduce the potential for broader clinical deployment. The table highlights both the progress made in this area and the need for more scalable, interpretable, and clinically transferable approaches to knee mobility assessment, especially vision-based frameworks that can preserve biomechanical relevance while supporting automated and real-world evaluation.



\section{Ankle: study details}
\label{supp:S3}
{ One line of work has focused on direct clinical assessment of ankle mobility and foot posture. Reilly et al. \cite{reilly2009role} studied the relationship between ankle dorsiflexion and foot posture in 20 patients with hip osteoarthritis (OA), 20 with knee OA and 20 healthy controls. Foot posture was assessed with the Foot Posture Index (FPI), and ankle dorsiflexion was measured clinically in a research setting. Patients with hip OA tended to have a more supinated foot posture and reduced dorsiflexion, whereas patients with knee OA had a more pronated foot posture and greater dorsiflexion than the hip OA group. Controls fell between these patterns. Simple clinical measures such as FPI and ankle dorsiflexion may therefore help distinguish forms of lower-limb OA and reveal compensatory foot and ankle behaviour \cite{reilly2009role}.}

{ A second group of studies used instrumented gait analysis to examine how ankle OA affects walking and how it changes after surgery. Valderrabano et al. \cite{valderrabano2007gait} studied patients with ankle OA before and after total ankle replacement (TAR), combining clinical scores with kinematic and kinetic gait assessment. A six-camera Vicon system, reflective markers and force plates were used, with follow-up from preoperative baseline to 12 months after surgery. Clinical status was assessed with the SF-36 and the American Orthopaedic Foot and Ankle Society (AOFAS) ankle score. TAR produced substantial improvement in both patient-reported outcomes and many gait variables. Several spatio-temporal parameters that were clearly abnormal before surgery were statistically comparable to healthy controls by 12 months. However, plantar flexion movement, medial ground reaction force, plantar flexion moment, adduction moment and ankle power still differed slightly from controls at final follow-up. TAR can therefore restore much of normal walking function, although some biomechanical deficits may persist after successful surgery \cite{valderrabano2007gait}.}

{ The wider impact of ankle OA has also been examined by comparing it with OA at other joints. Schmitt et al. \cite{schmitt2015effect} studied patients with isolated hip, knee and ankle OA, with 30 patients in each group and 15 healthy controls. An eight-camera motion capture system and four force plates recorded gait at self-selected speed. Isolated OA at any of the three joints altered walking mechanics not only at the affected joint but across the kinetic chain. For ankle OA this means ankle pathology should not be viewed in isolation, since it can influence movement at the knee and hip as well. This supports a more integrated biomechanical interpretation of gait \cite{schmitt2015effect}.}

{ Taken together, these} studies show that ankle-related mobility problems influence gait through changes in dorsiflexion, plantar flexion, foot posture, ground reaction forces, joint { moments} and whole-limb coordination \cite{reilly2009role,valderrabano2007gait,schmitt2015effect}. { Direct} clinical measures and { laboratory} gait analysis provide useful and complementary information. { The ankle literature remains small compared with the} hip and knee { literature, so ankle mobility is still} underexplored and there is considerable scope for more comprehensive, { scalable} and clinically applicable assessment.




{\scriptsize
\setlength{\tabcolsep}{2pt}
\renewcommand{\arraystretch}{1.03}
\begin{longtable}{p{1.9cm} p{1.6cm} p{1.9cm} p{1.9cm} p{1.7cm} p{1.9cm} p{2.0cm} p{3.0cm}}
\caption{Critical comparison of existing literature on ankle-related mobility disorders and gait analysis.}
\label{ankle_critical_extended_summary} \\
\hline
Literature & Cohort & Focus & Data / Setup & Analysis & Contribution & Strength & Limitation \\
\hline
\endfirsthead

\hline
Literature & Cohort & Focus & Data / Setup & Analysis & Contribution & Strength & Limitation \\
\hline
\endhead

Reilly et al. \cite{reilly2009role}
& 20 hip OA, 20 knee OA, 20 controls
& Foot posture and ankle dorsiflexion
& Clinical goniometric assessment; FPI scoring in research clinic
& Correlation and group comparison
& Showed that hip and knee OA are associated with different foot posture patterns and dorsiflexion ranges
& Clinically simple, low-cost, and directly relevant for screening lower-limb alignment differences
& Limited to static or clinic-based measures; does not capture full gait mechanics or dynamic ankle loading during walking \\
\hline

Valderrabano et al. \cite{valderrabano2007gait}
& 15 ankle OA, 15 matched controls
& Pre- and post-TAR gait recovery
& 6-camera Vicon system; 15 markers; 2 force plates; SF-36; AOFAS
& Repeated-measures ANOVA; t-tests; SPSS; kinematic and kinetic analysis
& Demonstrated substantial functional and biomechanical improvement after TAR, with most gait variables approaching control values by 12 months
& Strong longitudinal design; combines patient-reported outcomes with detailed gait kinetics and kinematics
& Small cohort and highly laboratory-based protocol; some biomechanical deficits remained at 12 months, suggesting incomplete restoration of normal ankle function \\
\hline

Schmitt et al. \cite{schmitt2015effect}
& 30 hip OA, 30 knee OA, 30 ankle OA, 15 controls
& Isolated joint OA and whole-limb walking mechanics
& 8-camera marker-based gait analysis; 4 force plates; 7 trials at self-selected speed
& Comparative kinematic and kinetic analysis
& Showed that isolated OA at the ankle alters walking mechanics not only locally but across the lower-limb kinetic chain
& Valuable comparative design across hip, knee, and ankle OA; supports integrated biomechanical interpretation
& Does not focus exclusively on ankle OA; broader whole-limb comparison reduces the depth of ankle-specific analysis \\
\hline

\end{longtable}
}

Table~\ref{ankle_critical_extended_summary} provides a critical comparison of the literature studied on ankle-related mobility disorders and gait analysis. Viewed together, these studies cover ankle osteoarthritis, lower-limb osteoarthritis more broadly, and post-surgical recovery, while examining issues such as mobility restriction, foot posture, surgical outcome, and whole-limb compensation. Methodologically, the literature combines direct clinical measures, such as goniometric assessment and the Foot Posture Index, with more detailed laboratory-based gait analysis using multi-camera systems, force plates, and patient-reported outcome scores. The comparison shows that ankle-related disorders affect gait in several interconnected ways, including changes in dorsiflexion, plantar flexion, foot posture, ground reaction forces, joint moments, and coordination across the lower-limb kinetic chain. One of the strengths of the current evidence base is that it brings together both direct clinical assessment and detailed biomechanical analysis, allowing ankle mobility to be understood from complementary functional and mechanistic perspectives. Also, it is evident that ankle-focused gait research remains relatively limited in both scale and depth when compared with the hip and knee literature. Most studies are based on small cohorts and highly controlled laboratory protocols, which can restrict broader clinical translation. Taken together, these observations underline the need for more scalable, interpretable, and clinically transferable frameworks for ankle mobility assessment.



\section{Human action recognition: methods and datasets}
\label{supp:S4}
Vemulapalli et al. \cite{vemulapalli2014human} proposed a mathematically grounded representation of 3D skeletons as points in a Lie group, specifically the Special Euclidean group SE(3). Their method captured the rotational and translational relationships among body parts and encoded the resulting trajectories into discriminative feature vectors before classification with an SVM. Evaluated on the MSR-Action3D, UTKinect-Action, and Florence3D-Action datasets, the approach consistently outperformed representations based on joint positions, pairwise relative positions, joint angles, and body-part locations. Similarly, Li et al. \cite{li2010action} proposed a depth-based action recognition approach using a bag of 3D points, showing that compact structural representations can achieve strong performance without dense motion modelling.

\begin{revision}
Boswell et al. \cite{boswell2021neural} estimated peak knee adduction moment using anatomical landmarks from motion capture and a reference moment calculated by inverse dynamics. The cohort comprised 86 participants, of whom 64 had knee osteoarthritis. The authors used two-dimensional projections of the landmarks to simulate camera views. Across all walking conditions in the held-out eight-person test set, $R^2$ for step-level peak moment was 0.85 using frontal-plane landmarks and 0.14 using sagittal-plane landmarks. The comparison supports further testing of frontal-view estimates using landmarks extracted from video.
\end{revision}

The importance of dataset relevance is clearly illustrated by Vafadar et al. \cite{vafadar2021novel}, who developed a gait-specific dataset to improve markerless motion capture during walking. Their results showed that fine-tuning a pose estimation model on gait-specific data substantially improved joint localisation accuracy, especially at the knee and ankle. This is a valuable finding, because it suggests that models trained on general-purpose action datasets may not be sufficient for clinical gait analysis unless they are adapted to task-specific movement patterns. Related work on video-based gait analysis and markerless movement assessment also supports the growing potential of pose-guided approaches for rehabilitation and clinical evaluation \cite{prakash2018vision,lv2022analysis,armstrong2022novel,stenum2024clinical}.

Other studies have focused more explicitly on feature representation. Caetano et al. \cite{caetano2019skeleton} introduced the Tree Structure Reference Joints Image (TSRJI), a skeleton-image representation designed to preserve joint relationships for CNN-based action recognition. Ke et al. \cite{ke2017skeletonnet} proposed SkeletonNet, which used carefully designed vector-based features to encode spatial body-part relationships before classification. Wang et al. \cite{wang2012mining} developed an action-recognition framework based on local occupancy patterns and a Fourier Temporal Pyramid (FTP), while Ahad et al. \cite{ahad2021action} proposed a kinematic posture feature set that combines angular and positional skeletal information. These studies are especially relevant for joint mobility analysis because they retain a closer connection between computational features and the underlying mechanics of movement.

{\scriptsize
\setlength{\tabcolsep}{3pt}
\renewcommand{\arraystretch}{1.08}
\begin{longtable}{m{0.7cm} m{2.0cm} m{0.9cm} m{1.5cm} m{2.0cm} m{2.3cm} m{2.1cm} m{2.0cm} m{0.7cm} m{0.7cm}}
\caption{HAR-based studies relevant to JMPs.}
\label{HAR} \\
\hline
\textbf{SN} & \textbf{Author} & \textbf{Year} & \textbf{Joint / Region} & \textbf{Problem Type} & \textbf{Method Used} & \textbf{Tools / Datasets} & \textbf{Metrics / Parameters} & \textbf{CV} & \textbf{HAR} \\
\hline
\endfirsthead

\hline
\textbf{SN} & \textbf{Author} & \textbf{Year} & \textbf{Joint / Region} & \textbf{Problem Type} & \textbf{Method Used} & \textbf{Tools / Datasets} & \textbf{Metrics / Parameters} & \textbf{CV} & \textbf{HAR} \\
\hline
\endhead

1 & W. Li et al. & 2010 & Full body & Action recognition & 3D silhouette representation with BMLD & Depth maps at 15 FPS; custom dataset & Accuracy: 93.4\%, 92.9\%, 96.3\% (AS1--AS3); average 94.2\% & Yes & Yes \\
\hline
2 & J. Wang et al. & 2012 & Lower limbs / full body & Actionlet-based action analysis & Fourier Temporal Pyramid + Actionlet Ensemble & MSR-Action3D, MSRDailyActivity3D, CMU MoCap & Accuracy: 88.2\%, 85.75\%, 98.13\% & Yes & Yes \\
\hline
3 & R. Vemulapalli et al. & 2014 & Full body & Action classification & Lie group representation + SVM & MSR-Action3D, UTKinect-Action, Florence3D-Action & Recognition rate: 92.46\%, 97.08\%, 90.88\% & Yes & Yes \\
\hline
4 & Y. Du et al. & 2015 & Legs and trunk & Sequence-based HAR & Hierarchical bidirectional RNN (HBRNN-L) & MSR Action3D, Berkeley MHAD, HDM05 & Mean accuracy: 94.49\%, 100\%, 96.92\% & Yes & Yes \\
\hline
5 & Q. Ke et al. & 2017 & Full body & Skeleton-based temporal pattern learning & SkeletonNet (CNN-based) & NTU RGB+D, SBU Kinect, CMU dataset & Accuracy: 81.16\% (cross-view), 93.47\%, 84.83\% & Yes & Yes \\
\hline
6 & C. Caetano et al. & 2019 & Full body & Skeleton-image action recognition & TSRJI + CNN & NTU RGB+D 60 and NTU RGB+D 120 & Accuracy: 73.3\% (cross-subject), 80.3\% (cross-view) & Yes & Yes \\
\hline
7 & T. Ahmad et al. & 2019 & Full skeleton & Attention-guided action recognition & Attention-Joint GCN & Stanford 40, PKU-MMD, NTU RGB+D & Mean accuracy: 84.6\%, 95.5\%, 90.7\% & Yes & Yes \\
\hline
8 & A. Viswakumar et al. & 2019 & Lower limbs / gait & Human gait analysis & OpenPose-based pose extraction & OpenPose; video-based gait analysis & Joint trajectories and gait pattern extraction & Yes & Yes \\
\hline
9 & A. Rohan et al. & 2020 & Foot and toe & Abnormal gait recognition & 2D pose estimation + CNN & PAF + VGG-19; custom gait videos & Accuracy: 97.3\%; sensitivity; precision; FPS & Yes & Yes \\
\hline
10 & M. A. Boswell et al. & 2021 & Knee & KAM prediction in OA & \rev{Neural network with projected motion-capture landmarks} & \rev{Simulated 2D views and inverse-dynamics reference} & \rev{Frontal $R^2=0.85$, sagittal $R^2=0.14$} & \rev{No} & \rev{No} \\
\hline
11 & S. Vafadar et al. & 2021 & Hip, knee, ankle & Joint localisation / angle estimation & Pose estimation with gait-specific fine-tuning & Human3.6M, ENSAM, VICON reference & Joint position error; Euclidean distance; Bland--Altman agreement & Yes & Yes \\
\hline
12 & M. A. R. Ahad et al. & 2021 & Knee and hip / full skeleton & HAR feature fusion & Kinematic Posture Features + SVM / CNN-RNN / ConvRNN & Five public datasets & Accuracy up to 98.44\%; KPF, LJPF, AJPF features & Yes & Yes \\
\hline
13 & A. Suhendar et al. & 2022 & Full body & Classifier benchmarking for HAR & OpenPose + SVM / ANN / LSTM / Transformer & KARD dataset & Accuracy, precision, recall, F1; best accuracy 99.4\% & Yes & Yes \\
\hline
14 & S. K. Yadav et al. & 2022 & Full body & Skeleton-based action recognition & ConvLSTM + guided feature learning & Custom dataset (130k samples) & Accuracy: 98.85\%; precision: 97.89\%; recall; F1-score & Yes & Yes \\
\hline
15 & X. Lv et al. & 2022 & Knee, ankle, foot & Arthritis-related gait recognition & Pose estimation + entropy-based feature fusion & COCO2017 + custom dataset & Accuracy: 97.19\% & Yes & Yes \\
\hline

\end{longtable}
}

Temporal modelling has also become increasingly important. Suhendar et al. \cite{suhendar2022skeletal} compared SVM, Artificial Neural Network (ANN), LSTM, and transformer models for skeleton-based HAR and found that the transformer achieved the strongest classification performance, particularly when more temporal frames were used. Yadav et al. \cite{yadav2022skeleton} showed that a ConvLSTM model, combining convolutional and recurrent learning, could outperform standard Convolution Neural Network (CNN) and LSTM baselines. Earlier, Du et al. \cite{du2015hierarchical} demonstrated the value of hierarchical bidirectional recurrent modelling for skeleton-based action recognition. Together, these studies underline the importance of temporal structure in human movement analysis, especially for gait, where clinically meaningful differences often emerge through the timing and coordination of motion rather than through isolated static postures.

Attention-based methods have added a further layer of refinement. Ahmad et al. \cite{ahmad2019action} proposed an attention-joint graph convolutional network that identifies the most relevant joints for different actions and then models their relationships explicitly. This is particularly promising for joint mobility applications, because it supports the selective emphasis of anatomically meaningful joints rather than treating all body points as equally informative. In a similar spirit, Uddin et al. \cite{uddin2011human} demonstrated that joint-angle-based features derived from a 3D body model can also support highly accurate HAR, especially when combined with dimensionality reduction and probabilistic sequence modelling.

Based on the above discussion, it is clear that HAR has evolved well beyond simple action recognition. It now incorporates a diverse set of techniques, including pose estimation, skeletal representation learning, temporal modelling, graph-based reasoning, and clinically relevant biomechanical prediction \cite{kong2022human,debnath2022review,hulleck2022present}. For joint mobility problems, the importance of this body of work lies in its shift from recognising what action is being performed to analysing how that movement is executed. In this setting, the most valuable contributions are not simply those achieving high recognition accuracy, but those that preserve biomechanical meaning, enable interpretability, and offer realistic potential for clinical application. Therefore, the reviewed studies establish a strong methodological basis for applying HAR to joint mobility problems. Table~\ref{HAR} presents a comparison of existing HAR-based literature relevant to joint mobility problems, highlighting the affected joints, problem focus, modelling strategy, computational setting, and reported evaluation metrics. The inclusion of performance measures makes it clear that the literature has progressed from general full-body action recognition toward more targeted gait- and joint-related assessment, including abnormal gait detection, joint localisation, and clinically meaningful biomechanical prediction. Additionally, Table~\ref{HAR} shows that reported metrics are not yet fully standardised across studies, which makes direct comparison more difficult and reinforces the need for more unified evaluation protocols in clinically oriented HAR research.







\section{Vision-based gait analysis: study details}
\label{supp:S5}
Lv et al. \cite{lv2022analysis} proposed a gait analysis framework based on AlphaPose combined with sample entropy fusion for the assessment of patients with knee osteoarthritis. Their method used a top-down pose estimation strategy and then transformed the extracted joint trajectories into one-dimensional signals for subsequent processing. Ensemble empirical mode decomposition was applied to handle non-stationary noise and reduce baseline drift, after which sample entropy was used to characterise differences in joint signal complexity between patient and control groups. The study found statistically significant entropy differences, particularly at the heel and ankle, suggesting that foot landing dynamics are especially informative for lower-limb arthritis. This work is valuable because it shows how pose estimation can be combined with signal-processing methods to derive clinically discriminative gait markers from simple video.

Ju et al. \cite{ju1996cardboard} provided one of the earliest influential frameworks for modelling articulated motion in image sequences. Their parameterised ``cardboard people'' model represented the human body as a set of connected planar segments and demonstrated how articulated motion could be captured efficiently despite occlusion and partial visibility. Although this work predates modern deep learning, it remains conceptually important because it established a structured representation of articulated motion that later computer vision and gait analysis methods would build upon.

Leightley et al. \cite{leightley2016automated} proposed a single-depth-sensor framework for recognising and quantifying human mobility. Their system grouped joints into anatomically meaningful regions and combined feature engineering with machine learning to estimate clinically relevant mobility measures. By computing variables such as Euler angles, Euclidean distances, centre of mass, and body lean angle, and then classifying mobility quality using support vector machines, the authors demonstrated that relatively low-cost sensing could still produce useful movement assessment. The study is particularly relevant because it moved beyond simple action recognition and aimed to provide interpretable mobility outcomes.

Fan et al. \cite{fan2020gaitpart} introduced GaitPart, a part-based framework for gait recognition that used focal convolution and temporal aggregation. Instead of treating gait as a single global pattern, the model extracted fine-grained features from horizontally partitioned body regions and captured short-range temporal dynamics within each part. Evaluated on the CASIA-B and OU-MVLP datasets, the method outperformed competing approaches and demonstrated the value of part-aware modelling for gait recognition. Although the primary focus was biometric gait recognition, the methodological insight is highly relevant to joint mobility assessment because clinically meaningful gait changes are often localised to specific body segments rather than distributed uniformly across the whole body.

Cotton et al. \cite{cotton2022transforming} proposed a transformer-based approach for estimating gait parameters from monocular video. Their model used 3D joint trajectories derived from PosePipe and mapped them, together with subject height, to temporal and kinematic gait features. The results showed strong agreement with laboratory ground truth for cadence, step time, double support time, gait velocity, step length, and lower-limb joint angles. This study is particularly significant because it demonstrates that monocular video, when combined with modern sequence modelling, can provide quantitatively meaningful gait parameters that \rev{warrant further validation for each intended clinical measurement.}

Prakash et al. \cite{prakash2018vision} proposed a markerless framework for identifying lower-limb joints, including the shoulder, hip, and knee, from conventional video. Their method relied on silhouette extraction, horizontal body segmentation, and geometric reasoning to identify limb configuration and estimate gait parameters such as step size, cadence, stride length, gait cycle time, and stance-swing proportions. The reported values were close to reference standards, illustrating that even relatively simple computer vision pipelines can recover useful gait measures when carefully designed.

Uccheddu et al. \cite{uccheddu2022rgb} presented an RGB-D approach for measuring hip and knee range of motion during home-based rehabilitation. By combining calibration, body-part segmentation, and depth-based angle estimation, the system achieved angular errors of fewer than five degrees and showed no significant difference from motion capture for most movements. This study is important because it extends markerless gait and movement analysis into rehabilitation contexts where ease of use and home deployment are essential.

Armstrong et al. \cite{armstrong2022novel} compared a markerless camera-based system with a marker-based knee kinesiography system to determine which better reflected natural walking biomechanics. Their findings suggested that markerless analysis produced a more representative walking pattern, particularly when compared with motion recorded while participants wore the KneeKG device. This is an important contribution, because it highlights not only the feasibility of markerless gait analysis, but also its potential advantage in reducing measurement artefacts introduced by instrumentation.

Juri et al. \cite{juri2020characterization} evaluated a low-cost OpenPose-based markerless gait analysis system against an inertial measurement unit reference. Although the angular trajectories of hip, knee, and ankle motion were broadly similar across systems, the markerless approach showed greater variability and some under- or overestimation in flexion and extension angles. Even so, the study demonstrated that clinically relevant sagittal plane trajectories can be recovered from a simple stereo video setup, which is promising for accessible gait assessment.

Ong et al. \cite{ong2017efficacy} directly compared a markerless system with a marker-based motion capture system for 3D joint-angle estimation during gait. Their results showed root mean square differences of less than one degree for most joint rotations, with significant differences appearing mainly at the ankle for inversion/eversion and internal/external rotation. This study is especially valuable because it provides strong quantitative evidence that carefully calibrated markerless video can approximate gold-standard gait measurements under controlled conditions.

The aforementioned literature highlights that computer vision and HAR have advanced substantially in their ability to analyse gait from ordinary video, evolving from early articulated motion models to more clinically relevant markerless gait assessment frameworks \cite{ju1996cardboard,prakash2018vision,lv2022analysis,cotton2022transforming}. A major strength of this research lies in its capacity to recover meaningful gait descriptors without relying on complex marker-based instrumentation, making movement analysis more practical, scalable, and accessible. At the same time, the existing studies also make it clear that important challenges remain, particularly in relation to accuracy, robustness to noise, and consistency of joint estimation, especially when these methods are transferred from controlled laboratory settings to real-world clinical environments \cite{ong2017efficacy,juri2020characterization,armstrong2022novel,stenum2024clinical}. Additionally, the reviewed studies provide a strong methodological foundation for scalable and interpretable video-based gait analysis, with clear potential to enhance joint mobility assessment in healthcare and rehabilitation.

{\scriptsize
\setlength{\tabcolsep}{3pt}
\renewcommand{\arraystretch}{1.08}
\begin{longtable}{p{1.8cm} p{0.8cm} p{1.4cm} p{1.8cm} p{2.0cm} p{2.2cm} p{2.2cm} p{3.3cm}}
\caption{Critical comparison of computer vision and HAR-based studies for gait analysis from 2D or markerless video.}
\label{cv_har_gait_critical} \\
\hline
\textbf{Reference} & \textbf{Year} & \textbf{Type} & \textbf{Input / Setup} & \textbf{Key Objective} & \textbf{Core Strength} & \textbf{Key Limitation} & \textbf{Critical Contribution to JMP} \\
\hline
\endfirsthead

\hline
\textbf{Reference} & \textbf{Year} & \textbf{Type} & \textbf{Input / Setup} & \textbf{Key Objective} & \textbf{Core Strength} & \textbf{Key Limitation} & \textbf{Critical Contribution to JMP} \\
\hline
\endhead

Ju et al. \cite{ju1996cardboard}
& 1996
& Foundational articulated motion model
& 2D image sequences; parameterised body segments
& Represent articulated human motion efficiently
& Early structured representation of body motion; robust to partial occlusion
& Not designed for clinical gait analysis; lacks direct biomechanical outputs
& Provides an important conceptual basis for later markerless motion and pose-based gait analysis frameworks \\
\hline

Leightley et al. \cite{leightley2016automated}
& 2016
& Depth-based mobility assessment
& Single depth sensor; joint-group MoCap features
& Recognise motion and quantify mobility
& Moves beyond recognition toward clinically meaningful mobility scoring
& Uses depth sensing rather than ordinary RGB video; mobility labels are relatively coarse
& Shows that low-cost sensing can support interpretable mobility assessment, but with limited joint-specific precision \\
\hline

Ong et al. \cite{ong2017efficacy}
& 2017
& Markerless vs marker-based validation
& Two cameras + marker-based optical system
& Evaluate 3D joint-angle accuracy of markerless gait analysis
& Strong direct comparison against gold-standard motion capture
& Small sample of healthy young males only; controlled treadmill setting
& Demonstrates that markerless gait analysis can approximate laboratory-grade kinematics under controlled conditions \\
\hline

Prakash et al. \cite{prakash2018vision}
& 2018
& Markerless silhouette-based gait analysis
& Single RGB camera; silhouette segmentation
& Extract gait parameters and lower-limb joint locations
& Simple and low-cost pipeline; clinically meaningful spatio-temporal outputs
& Relies on handcrafted rules and silhouette quality; limited robustness in unconstrained settings
& Illustrates that useful gait descriptors can be recovered from conventional video without markers \\
\hline

Fan et al. \cite{fan2020gaitpart}
& 2020
& Part-based gait recognition
& Silhouette video; CASIA-B and OU-MVLP
& Improve gait recognition using local body-part modelling
& Strong part-aware spatial-temporal representation; excellent benchmark performance
& Focuses on biometric recognition rather than clinical gait interpretation
& Methodologically valuable because local body-part modelling is relevant to joint-specific movement abnormalities \\
\hline

Juri et al. \cite{juri2020characterization}
& 2020
& Low-cost markerless clinical gait analysis
& Two webcams + OpenPose + IMU reference
& Assess feasibility of OpenPose-based gait analysis
& Strong practical focus; compares markerless estimates with wearable reference
& Single healthy participant; variability remains high; limited generalisability
& Highlights both the promise and current instability of low-cost markerless gait assessment in clinical use \\
\hline

Lv et al. \cite{lv2022analysis}
& 2022
& Pose-based arthritis discrimination
& 2D video; AlphaPose + EEMD + sample entropy
& Distinguish knee OA gait from healthy gait
& Combines pose estimation with signal analysis; identifies clinically meaningful entropy differences
& Small cohort; statistical discrimination rather than full predictive modelling
& Shows that pose-derived temporal complexity measures can reveal disease-relevant gait alterations \\
\hline

Cotton et al. \cite{cotton2022transforming}
& 2022
& Monocular gait parameter estimation
& Monocular video + PosePipe + transformer
& Predict temporal and kinematic gait parameters
& Large laboratory dataset; strong agreement with ground truth across several gait variables
& Still trained and tested in controlled lab settings; some parameters less accurate than others
& One of the strongest examples of translating monocular video into clinically meaningful gait parameters \\
\hline

Uccheddu et al. \cite{uccheddu2022rgb}
& 2022
& RGB-D rehabilitation assessment
& RGB-D cameras; healthy and impaired subjects
& Measure hip and knee ROM in home rehabilitation
& Good practical relevance; direct comparison with motion capture; low angular error
& Focuses on ROM rather than full gait biomechanics; requires RGB-D rather than ordinary RGB
& Important for rehabilitation monitoring, especially where joint ROM is the primary outcome \\
\hline

Armstrong et al. \cite{armstrong2022novel}
& 2022
& Markerless vs instrumented system comparison
& Azure Kinect + KneeKG
& Compare natural gait representation across systems
& Strong translational value; questions whether instrumented systems alter natural gait
& Narrow focus on a few kinematic measures; limited cohort and treadmill context
& Suggests that markerless systems may better preserve natural walking behaviour than instrumented clinical systems \\
\hline
\end{longtable}
}

Table~\ref{cv_har_gait_critical} presents a critical comparison of literature that use computer vision and HAR for gait analysis from 2D or markerless video. The comparison shows a clear methodological progression in the field. Earlier work was primarily concerned with representing articulated human motion in a computationally efficient way, whereas more recent studies have focused on extracting clinically meaningful gait parameters, validating markerless systems against reference technologies, and distinguishing pathological from healthy movement patterns. It also highlights that the literature now spans a broad spectrum of objectives, including motion representation, gait recognition, disease discrimination, rehabilitation monitoring, and quantitative parameter estimation. A major strength that emerges from this comparison is the growing ability of video-based methods to recover meaningful gait descriptors without relying on complex marker-based instrumentation. This makes the approach attractive for scalable, low-cost, and potentially more natural movement assessment. It also shows that the evidence base is still quite heterogeneous. Differences in sensing modality, cohort size, evaluation protocol, and clinical target make direct comparison difficult, while several studies remain limited by small samples, controlled laboratory settings, or narrow validation scenarios. In general, the reviewed literature provide a strong methodological foundation for video-based joint mobility assessment, but they also make clear that wider clinical validation, improved robustness, and more standardised evaluation frameworks are still needed before these methods can be adopted more confidently in routine practice.





\section{Human pose estimation: study details}
\label{supp:S6}
{\scriptsize
\setlength{\tabcolsep}{3pt}
\renewcommand{\arraystretch}{1.08}
\begin{center}
\begin{longtable}{m{0.7cm} m{2.2cm} m{2.2cm} m{2.0cm} m{2.1cm} m{3.0cm} m{2.8cm}}
\caption{Summary of human pose estimation (HPE) applications relevant to posture recognition, gait analysis, rehabilitation, and joint mobility assessment.}
\label{HPE} \\
\hline
\textbf{S.No.} & \textbf{Author (Year)} & \textbf{Application Area} & \textbf{Pose Framework} & \textbf{Dataset} & \textbf{Methodology} & \textbf{Metrics} \\
\hline
\endfirsthead

\hline
\textbf{S.No.} & \textbf{Author (Year)} & \textbf{Application Area} & \textbf{Pose Framework} & \textbf{Dataset} & \textbf{Methodology} & \textbf{Metrics} \\
\hline
\endhead

1 & Ma et al. (2022) \cite{ma2022wushu} 
& Wushu posture recognition 
& MediaPipe 
& 400 images per pose 
& Rule-based pose comparison using Euclidean distance, vector relations, and cosine similarity 
& Recognition accuracy discussed qualitatively; performance constrained by single-camera setup \\ \hline

2 & Li et al. (2020) \cite{li2020human} 
& Lower-limb rehabilitation assessment 
& HRNet / MobileNet 
& COCO2017 + custom rehabilitation dataset 
& Lightweight HRNet with depthwise and dilated convolutions for efficient pose estimation 
& AP = 70.8 on COCO; AP = 88.3 on custom dataset; model size reduced substantially relative to HRNet \\ \hline

3 & Sharma et al. (2022) \cite{sharma2022surya} 
& Yoga pose recognition and correction 
& MediaPipe + LSTM 
& Webcam-based Surya Namaskar videos 
& Pose landmarks extracted by MediaPipe and classified using temporal LSTM modelling 
& Overall accuracy = 96.67\% \\ \hline

4 & Viswakumar et al. (2019) \cite{viswakumar2019human} 
& Markerless gait analysis under varying lighting and clothing 
& OpenPose (BODY\_25) 
& 10 healthy participants 
& 2D keypoint extraction with joint-angle computation under different illumination and clothing conditions 
& Knee-angle error = 13.34\%--16.78\% for lighting variation; 17.8\%--18.29\% for clothing variation \\ \hline

5 & Stenum et al. (2021) \cite{stenum2021two} 
& Two-dimensional video-based gait analysis 
& OpenPose 
& \rev{31 healthy participants analysed} 
& \rev{Sagittal views of overground walking compared with synchronised 3D motion capture} 
& \rev{Mean absolute differences of 0.02 s for individual step times and 0.049 m for individual step lengths} \\ \hline

6 & Van Hooren et al. (2023) \cite{van2023accuracy} 
& Running kinematics 
& DeepLabCut / OpenPose 
& In-laboratory running videos 
& Markerless kinematic estimation compared with marker-based motion capture 
& OpenPose showed better agreement after alignment; accuracy varied across speed and conditions \\ \hline

7 & Sheu et al. (2023) \cite{sheu2023improvement} 
& Real-time pose estimation and action recognition 
& IFC Network 
& Penn Action, Sub-JHMDB, UCF101 
& Intensive Feature Consistency network for inter-frame pose refinement and temporal consistency 
& PCK-body = 99.1\%; PCK-torso = 94.7\%; real-time speed = 31 FPS \\ \hline

8 & Marusic et al. (2023) \cite{marusic2023evaluating} 
& Rehabilitation exercise monitoring 
& OpenPose / BlazePose / Kinect 
& Keraal rehabilitation dataset 
& Motion-sequence classification using pose estimation and Gaussian Mixture Models 
& OpenPose and BlazePose achieved performance comparable to Kinect \\ \hline

9 & Blasco-García et al. (2024) \cite{blasco2024computer} 
& Fall-risk evaluation in older adults 
& RGB-D pose estimation 
& Laboratory and nursing-home assessments 
& Extraction of gait and balance parameters from clinical tests such as TUG and Tinetti 
& Effective for remote monitoring and longitudinal fall-risk assessment; performance reported as clinically useful \\ \hline

10 & Wang et al. (2024) \cite{wang2024markerless} 
& Markerless gait analysis 
& HuMoR 
& 14 healthy participants 
& Single-camera 3D keypoint estimation with comparison against marker-based motion tracking 
& Good to strong correlations for major spatiotemporal parameters; weaker results for some ankle and pelvic measures \\ \hline

11 & Stenum et al. (2024) \cite{stenum2024clinical} 
& Clinical gait analysis across patient populations 
& OpenPose 
& Healthy adults, post-stroke patients, and Parkinson's disease patients 
& Multi-view video-based gait analysis validated against 3D motion capture 
& Strong correlations for key gait parameters, with $r \geq 0.872$ \\ \hline

12 & Doong (2025) \cite{doong2025predicting} 
& Workplace postural risk assessment 
& 3DMPPE / 3DWild 
& Multi-source REBA-related image datasets 
& Weakly supervised learning with domain adaptation for ergonomic risk classification 
& Best balanced accuracy = 0.761; best F1 score = 0.748 \\ \hline

13 & Yan et al. (2025) \cite{yan2025efficient} 
& Range of motion assessment on edge devices 
& Reduced Quantized Hourglass (RQ-HG) 
& Custom ROM video dataset 
& Quantized monocular ROM assessment model deployed on Raspberry Pi 
& Joint prediction accuracy = 93.6\%; RMSE = 3.21$^\circ$--4.25$^\circ$ \\ \hline

14 & Schoenwether et al. (2025) \cite{schoenwether2025reliability} 
& Gait biomechanics reliability analysis 
& KinaTrax 
& 9 healthy adults across repeated sessions 
& Markerless gait capture assessed using within- and between-session reliability analysis 
& ICC $>$ 0.9 for most spatiotemporal parameters; strongest reliability in sagittal-plane kinematics \\ \hline

\end{longtable}
\end{center}
}

A number of studies have concentrated on posture recognition using lightweight and publicly accessible frameworks. For instance, Ma et al. \cite{ma2022wushu} developed a Wushu posture-recognition system using MediaPipe and a standard camera. Their framework used the head, hand, and body APIs of MediaPipe to identify seven standard Wushu movements, including fist, palm, hook, tiger talon, forward lunge, horse stance, and empty step. The recognition logic was formulated using Euclidean distance, vector relationships, and cosine similarity between extracted landmarks and pre-defined posture templates. The resulting system was implemented in Python to assist users in assessing Wushu postures without requiring expert supervision. Although the study demonstrated the feasibility of low-cost posture analysis, the authors also acknowledged that the use of a single camera limited recognition precision, particularly for postures involving depth-dependent limb configurations. A similar lightweight strategy was adopted by Sharma et al. \cite{sharma2022surya}, who proposed a real-time system for recognising and correcting Surya Namaskar yoga postures. Their method combined MediaPipe-based skeletal landmark extraction with a LSTM network for temporal classification and achieved an overall accuracy of 96.67\%. Importantly, the system also provided corrective feedback by comparing the user’s pose to a reference posture, demonstrating that HPE can support not only recognition but also interactive movement guidance in smart-healthcare applications.

Pose estimation has also become increasingly important in rehabilitation, where efficiency, portability, and ease of deployment are critical. Li et al. \cite{li2020human} proposed a lightweight version of the High-Resolution Network (HRNet), enhanced with depthwise and dilated convolutions, for lower-limb rehabilitation exercises. Their dataset was collected using a nine-camera system with reflective markers and included 31 rehabilitation exercises performed according to physiotherapist-defined standards. The lightweight model achieved an average precision of 70.8 on COCO2017 and 88.3 on the custom rehabilitation dataset, while substantially reducing model size and computational complexity relative to the original HRNet. This is significant because it shows that clinically relevant pose-estimation models can be made more portable without a major sacrifice in performance, which is essential for future mobile-health and home-based rehabilitation systems.

A substantial body of work has used OpenPose for markerless gait analysis and related clinical tasks. Viswakumar et al. \cite{viswakumar2019human} investigated the performance of OpenPose for gait analysis under different lighting and clothing conditions. Using the BODY\_25 model, they recorded videos of ten healthy participants and examined the effect of dim versus bright illumination, as well as the influence of traditional clothing. Knee angles were estimated from thigh and shank vectors, and the resulting errors ranged from 13.34\% to 18.29\%, depending on the condition. Notably, OpenPose performed better than an MS Kinect-based depth system when participants wore dhotis, suggesting that some RGB-based pose estimators may be more robust than depth-based systems under challenging clothing conditions. \rev{Stenum et al. \cite{stenum2021two} compared OpenPose estimates from sagittal video with simultaneous 3D motion capture during overground walking. They analysed 31 of the 32 healthy participants in the dataset, excluding one recording from a different walking protocol. Agreement varied by measure and improved for some spatiotemporal variables when values were averaged within participants. The 2024 study discussed below examined clinical cohorts.}

The clinical scope of HPE-based gait analysis has broadened further in recent work. Stenum et al. \cite{stenum2024clinical} introduced a low-cost video-based pose estimation framework for clinical gait analysis across healthy adults, post-stroke patients, and individuals with Parkinson’s disease. By combining frontal and sagittal video views, the system achieved strong correlations with a gold-standard 3D motion capture system for key gait variables such as step time, step length, and gait speed. The ability to capture within-participant gait changes across diverse populations is especially important because it moves the literature closer to clinically transferable markerless assessment. Similarly, Blasco-García et al. \cite{blasco2024computer} proposed an RGB-D-based system for fall-risk assessment in older adults using pose landmarks extracted during clinical tests such as the Timed Up and Go (TUG) and Tinetti assessments. Their system provided objective data for balance and gait assessment without relying on wearable sensors, which \rev{motivates evaluation for remote or longitudinal monitoring in elderly care}. These studies show that HPE has progressed well beyond laboratory gait description and is increasingly being integrated into clinically relevant assessment workflows.

Other studies have examined the quantitative validity of markerless gait analysis relative to gold-standard systems. Ong et al. \cite{ong2017efficacy} evaluated a markerless video-based tracking system against marker-based motion capture and found that the root mean square differences between the two systems were generally below one degree for many joint rotations, although significant differences remained for some ankle motions. Juri et al. \cite{juri2020characterization} assessed a low-cost OpenPose-based 3D gait-analysis system against inertial measurement units and reported that angular trajectories were broadly similar, although OpenPose showed greater variability and underestimation in some flexion and extension measures. These findings are important because they reveal both the promise and the limitations of low-cost markerless systems: they can approximate kinematic trajectories reasonably well, but their stability and precision still depend strongly on calibration quality, joint type, and recording conditions. Wang et al. \cite{wang2024markerless} further showed that a single-camera HuMoR-based system could estimate spatiotemporal gait parameters and lower-limb joint angles with good to strong correlations relative to marker-based motion capture, although the results were clearly influenced by camera viewpoint and distance. This suggests that markerless gait analysis is becoming increasingly viable, but its robustness under uncontrolled real-world settings remains a major challenge.

Several recent studies have also aimed to improve the technical reliability and efficiency of HPE for real-time or edge-based applications. Sheu et al. \cite{sheu2023improvement} introduced the Intensive Feature Consistency (IFC) Network, which refines pose predictions across frames by enforcing both global and local feature consistency. Their framework achieved 99.1\% PCK-body and 94.7\% PCK-torso accuracy at 31 FPS, indicating that real-time refinement can substantially improve pose stability and accuracy. Yan et al. \cite{yan2025efficient} proposed a reduced and quantized hourglass model for range-of-motion estimation from a monocular RGB camera, specifically optimised for deployment on Raspberry Pi devices. Their system preserved 93.6\% joint-prediction accuracy while greatly reducing model size and inference cost, and achieved root mean squared errors between 3.21$^\circ$ and 4.25$^\circ$ relative to goniometer measurements. This represents an important step toward practical deployment in remote or resource-constrained clinical settings. Likewise, Schoenwether et al. \cite{schoenwether2025reliability} evaluated the KinaTrax markerless motion capture system and reported excellent within- and between-session reliability for most spatiotemporal gait parameters, particularly in the sagittal plane. Their results strengthen the argument that markerless systems can support repeated clinical assessment, although limitations remain for frontal- and transverse-plane kinematics.

The literature also includes studies that extend HPE beyond standard gait analysis into broader movement-assessment scenarios. Marusic et al. \cite{marusic2023evaluating} compared Kinect, OpenPose, and BlazePose for classifying rehabilitation exercises in a low-back-pain setting and found that OpenPose and BlazePose achieved accuracy comparable to Kinect. This is encouraging because it suggests that widely accessible RGB-camera systems may be sufficient for home-based rehabilitation monitoring when combined with suitable classifiers. Doong \cite{doong2025predicting} extended HPE into workplace ergonomics by using pose-derived joint angles and machine learning to predict postural risk levels through the Rapid Entire Body Assessment framework. By incorporating weakly supervised learning and multi-source domain adaptation, the study showed that pose estimation can support practical ergonomic risk stratification beyond conventional gait contexts. These applications demonstrate that HPE-based analysis is increasingly being treated not simply as a vision problem, but as a clinically and functionally meaningful decision-support tool.

In summary, this paper demonstrates that publicly available pose estimation frameworks have become versatile tools for posture recognition, gait analysis, rehabilitation support, and clinical monitoring. Their principal strength lies in enabling movement assessment from low-cost and accessible video data, often without the need for markers, force plates, or specialised laboratory infrastructure. At the same time, the literature also makes clear that performance remains sensitive to camera viewpoint, lighting, clothing, motion complexity, and the specific biomechanical variable being measured. Accuracy is generally strongest for large-joint and sagittal-plane motion, whereas more subtle multi-planar movement remains more difficult to estimate robustly \cite{ong2017efficacy,juri2020characterization,van2023accuracy,stenum2024clinical}. Consequently, although current HPE frameworks provide a strong methodological basis for scalable and interpretable movement analysis, further work is still needed to improve standardisation, generalisability, and clinical-grade robustness before these systems can be adopted routinely for joint mobility assessment in real-world healthcare settings.



\section{Extended method and benchmark comparisons}
\begin{revision}
These tables preserve the study-level methods and results that previously appeared in the main paper. Results refer to the stated tasks and datasets. Matching dataset names do not establish identical splits, preprocessing or evaluation settings, and action-recognition accuracy does not establish diagnostic performance for hip, knee or ankle disorders.
\end{revision}
{\scriptsize\setlength{\tabcolsep}{3pt}\begin{longtable}{c p{2.7cm} p{2.8cm} p{3.3cm} p{3.4cm} p{3.6cm}}
\caption{Chronological summary of relevant HAR-based studies related to gait, posture, mobility, and JMP.}
\label{HAR_M_A}\\
\hline
\textbf{SN} & \textbf{Author (Year)} & \textbf{Application} & \textbf{Methodology} & \textbf{Data Source/Literature Review} & \textbf{Key Results / Contribution} \\
\hline
\endfirsthead

\hline
\textbf{SN} & \textbf{Author (Year)} & \textbf{Application} & \textbf{Methodology} & \textbf{Data Source/Literature Review} & \textbf{Key Results / Contribution} \\
\hline
\endhead

1 & Shanon X. Ju \cite{ju1996cardboard} (1996) & Articulated human motion modelling & Parameterised articulated motion model for image sequences & Self-recorded image sequences & Early foundational framework for tracking articulated human motion in video. \\
\hline

2 & Zhengyou Zhang \cite{li2010action} (2010) & Action recognition using 3D silhouettes & Bag of 3D points, BMLD decoding, and NERF clustering & Self-recorded depth maps at 15 FPS & Achieved 94.2\% average classification accuracy. \\
\hline

3 & M. Z. Uddin \cite{uddin2011human} (2011) & HAR using body joint-angle features & Stereo vision, LDA, HMM, and joint-angle feature vectors & Self-recorded videos for 8 actions & Reported 98.13\% mean classification accuracy. \\
\hline

4 & J. Wang \cite{wang2012mining} (2012) & Action recognition with depth cameras & Local occupancy patterns, Fourier temporal pyramid, and actionlet ensemble & CMU MoCap, MSR-Action3D, and MSRDailyActivity3D & Achieved 88.2\%, 85.75\%, and 98.13\% accuracy, respectively. \\
\hline

5 & Raviteja Vemulapalli \cite{vemulapalli2014human} (2014) & 3D skeleton-based action recognition & Lie group representation of skeletons with DTW and SVM & MSR-Action3D, UTKinect-Action, and Florence3D-Action & Recognition rates of 92.46\%, 97.08\%, and 90.88\%, respectively. \\
\hline

6 & Y. Du \cite{du2015hierarchical} (2015) & Skeleton-based action recognition & Hierarchical bidirectional RNN with LSTM and BPTT & MSR Action3D, Berkeley MHAD, and HDM05 & Achieved 94.49\%, 100\%, and 96.92\% accuracy, respectively. \\
\hline

7 & D. Leightley \cite{leightley2016automated} (2016) & Mobility quantification and motion assessment & Depth-sensor motion capture, clustering, and SVM-based analysis & K3Da dataset & Achieved 87.1\% classification accuracy and demonstrated clinically oriented mobility quantification. \\
\hline

8 & Q. Ke \cite{ke2017skeletonnet} (2017) & Skeleton-based HAR & SkeletonNet with CNN-based deep part-feature learning & NTU RGB+D, SBU Kinect Interaction, and CMU MoCap & Reported 81.16\%, 93.47\%, and 84.83\% accuracy, respectively. \\
\hline

9 & A. Ong \cite{ong2017efficacy} (2017) & Markerless gait-analysis validation & Video-based markerless tracking compared with marker-based motion capture & 10 healthy participants & Showed low RMS differences for many joint rotations, supporting feasibility of markerless gait analysis. \\
\hline

10 & Chandra Prakash \cite{prakash2018vision} (2018) & Markerless gait parameter extraction & Vision-based silhouette processing with linear regression and geometric analysis & Self-recorded videos at 29 FPS & Produced gait parameters close to ground-truth and standard reference values. \\
\hline

11 & T. Ahmad \cite{ahmad2019action} (2019) & Skeleton-based action recognition using graph learning & Attention-joint graph CNN & Stanford 40 Actions, PKU-MMD, and NTU-RGBD & Achieved 84.6\%, 95.5\%, and 90.7\% accuracy, respectively. \\
\hline

12 & C. Caetano \cite{caetano2019skeleton} (2019) & Skeleton image representation for HAR & Tree Structure Reference Joints Image (TSRJI) with CNN & NTU RGB+D 60 and NTU RGB+D 120 & Achieved 73.3\% cross-subject and 80.3\% cross-view accuracy. \\
\hline

13 & Aditya Viswakumar \cite{viswakumar2019human} (2019) & Markerless gait analysis under varying conditions & OpenPose-based 2D pose estimation for gait-angle measurement & Self-recorded videos at 30 FPS & Demonstrated the effect of lighting and clothing on markerless gait estimation accuracy. \\
\hline

14 & Ali Rohan \cite{rohan2020human} (2020) & Abnormal gait detection & Pose estimation combined with VGG-19-based CNN classification & Self-recorded videos & Achieved 97.3\% classification accuracy for abnormal gait categories. \\
\hline

15 & Ying Li \cite{li2020human} (2020) & Lower-limb rehabilitation monitoring & Lightweight HRNet and MobileNet for pose estimation and exercise analysis & COCO2017 and custom rehabilitation dataset & Achieved AP of 70.8 on COCO and 88.3 on the custom rehabilitation dataset with reduced model complexity. \\
\hline

16 & Chao Fan \cite{fan2020gaitpart} (2020) & Gait recognition & Focal convolution, horizontal pooling, and temporal part aggregation & CASIA-B and OU-MVLP & Reported strong recognition performance across large-scale gait datasets. \\
\hline

17 & T. Juri \cite{juri2020characterization} (2020) & Low-cost 3D markerless gait analysis & OpenPose, linear triangulation, and IMU-based comparison & Single healthy participant & Reported ROM errors of 1.1$^\circ$ for hip and 3.8$^\circ$ for knee, supporting proof-of-concept use. \\
\hline

18 & Yucheng Chen \cite{chen2020monocular} (2020) & Monocular pose-estimation review & Survey of deep-learning-based monocular HPE methods & Literature review & Provides methodological background for 2D/3D pose estimation from monocular video. \\
\hline

19 & M. A. Boswell \cite{boswell2021neural} (2021) & \rev{Peak knee adduction moment estimation} & \rev{Neural network with projected gait landmarks} & \rev{Motion capture with simulated camera views} & \rev{For step-level peak moment, frontal $R^2=0.85$ and sagittal $R^2=0.14$.} \\
\hline

20 & S. Vafadar \cite{vafadar2021novel} (2021) & Markerless gait pose estimation improvement & Deep-learning-based markerless motion capture refined with gait-specific training & Human3.6M and ENSAM datasets & Reduced joint-estimation errors after fine-tuning on a gait-specific dataset. \\
\hline

21 & M. A. R. Ahad \cite{ahad2021action} (2021) & Action recognition from 3D skeleton joints & Kinematics posture features with SVM, RNN, ConvRNN, and CNNRNN & UTKinect-Action3D, KARD, MSR 3D Action Pairs, Florence 3D, and OAD & Reported consistently high accuracy across five public datasets. \\
\hline

22 & J. Stenum \cite{stenum2021two} (2021) & 2D video-based gait analysis & OpenPose-based gait measurement with comparison to motion capture & \rev{Overground walking in 31 healthy participants} & \rev{Agreement depended on the variable and whether estimates were averaged across steps.} \\
\hline

23 & Wang et al. \cite{wang2021measurement} (2021) & Lower-limb ROM estimation & Machine-vision-based joint ROM measurement through HMI & Experimental movement dataset & Demonstrated reliable lower-limb joint ROM estimation using vision-based analysis. \\
\hline

24 & Thomas Hellsten \cite{hellsten2021potential} (2021) & Rehabilitation-focused markerless motion analysis & Review of computer-vision-based markerless rehabilitation analysis & Literature review & Summarises the potential and limitations of markerless motion analysis for rehabilitation. \\
\hline

25 & A. Suhendar \cite{suhendar2022skeletal} (2022) & Benchmarking skeleton-based HAR classifiers & OpenPose with SVM, ANN, LSTM, and Transformer & KARD dataset & Transformer achieved 99.4\% accuracy, while SVM had the shortest training time. \\
\hline

26 & S. K. Yadav \cite{yadav2022skeleton} (2022) & Skeleton-based HAR using ConvLSTM & ConvLSTM, LSTM, and CNN with Kinect-based features & Custom dataset with 130,000 samples from 20 subjects & Achieved 98.85\% accuracy and 97.89\% precision. \\
\hline

27 & Xinyu Lv \cite{lv2022analysis} (2022) & Knee arthritis gait analysis & AlphaPose with TensorRT, EEMD, and sample entropy analysis & COCO2017 and self-recorded videos & Demonstrated significant entropy differences between patients and controls and improved processing speed. \\
\hline

28 & R. James Cotton \cite{cotton2022transforming} (2022) & Video-based spatiotemporal gait analysis & PosePipe with transformer-based gait-feature prediction & 7,331 frontal gait videos at 30 FPS & Reported strong correlations for multiple temporal, spatiotemporal, and kinematic gait parameters. \\
\hline

29 & F. Uccheddu \cite{uccheddu2022rgb} (2022) & Hip and knee ROM measurement in rehabilitation & RGB-D sensing with deep-learning-based segmentation & RGBD-HuDaAct and self-collected videos & Achieved average error below 5$^\circ$ for joint-angle estimation. \\
\hline

30 & K. Armstrong \cite{armstrong2022novel} (2022) & Markerless motion analysis for knee OA & Azure Kinect markerless motion capture compared with KneeKG & Pilot data / experimental setup & Showed statistically significant sensitivity to biomechanical differences, supporting low-cost clinical use. \\
\hline

31 & Jiewei Ma \cite{ma2022wushu} (2022) & Wushu posture recognition & MediaPipe with geometric pose comparison rules & Self-recorded posture images/videos & Demonstrated practical posture recognition using a low-cost single-camera setup. \\
\hline

32 & A. Sharma \cite{sharma2022surya} (2022) & Yoga posture recognition and correction & MediaPipe with LSTM-based temporal classification & Self-recorded videos & Achieved 96.67\% accuracy for Surya Namaskar posture recognition. \\
\hline

33 & Abdul Saboor \cite{saboor2020latest} (2020) & Review of gait analysis using wearable sensors and ML & Systematic review & Literature review & Provides useful methodological context for gait-analysis pipelines and machine-learning-based movement assessment. \\
\hline

34 & Alireza Sepas-Moghaddam \cite{sepas2022deep} (2022) & Deep gait recognition review & Survey of deep-learning-based gait-recognition methods & Literature review & Summarises modern deep gait-recognition architectures, datasets, and evaluation issues. \\
\hline

35 & Yu Kong \cite{kong2022human} (2022) & HAR review & Survey of action recognition and prediction methods & Literature review & Offers broad methodological context for HAR models relevant to movement analysis. \\
\hline

36 & Bappaditya Debnath \cite{debnath2022review} (2022) & Computer-vision-based rehabilitation assessment review & Review of vision-based rehabilitation and physical assessment methods & Literature review & Summarises rehabilitation-focused computer-vision methods relevant to joint mobility assessment. \\
\hline

37 & Bas van Hooren \cite{van2023accuracy} (2023) & Running kinematics with markerless vision systems & Comparison of OpenPose and DeepLabCut against marker-based motion capture & Running kinematics dataset & Showed that markerless methods are promising for field-scale data collection but remain variable at individual level. \\
\hline

38 & Ming-Hwa Sheu \cite{sheu2023improvement} (2023) & Improved pose estimation for action analysis & Intensive Feature Consistency network for pose refinement & Penn Action, Sub-JHMDB, and UCF101 & Achieved 99.1\% PCK-body and 94.7\% PCK-torso at real-time speed. \\
\hline

39 & Aleksa Marusic \cite{marusic2023evaluating} (2023) & Rehabilitation exercise assessment & Comparison of Kinect, OpenPose, and BlazePose with GMM-based classification & Keraal dataset & Showed that OpenPose and BlazePose can achieve accuracy comparable to Kinect in rehabilitation tasks. \\
\hline

40 & Hanwen Wang \cite{wang2024markerless} (2024) & Single-camera markerless gait analysis & HuMoR-based 3D body keypoint estimation from monocular video & 14 healthy participants & Reported good-to-strong correlations for several spatiotemporal and joint-angle measures, though viewpoint sensitivity remained important. \\
\hline

41 & Jesús Damián Blasco-García \cite{blasco2024computer} (2024) & Fall-risk assessment using pose estimation & RGB-D-based pose landmark extraction during clinical tests & Laboratory and nursing-home trials & Demonstrated clinically useful gait and balance assessment without wearable sensors. \\
\hline

42 & J. Stenum \cite{stenum2024clinical} (2024) & Clinical gait analysis using video-based pose estimation & Multi-view OpenPose-based clinical gait framework & Healthy adults, post-stroke, and Parkinson’s disease populations & Showed strong agreement with motion capture and good sensitivity to within-participant gait changes. \\
\hline

43 & Shing Hwang Doong \cite{doong2025predicting} (2025) & Postural risk prediction & HPE with multi-source domain adaptation and weakly supervised learning & Multi-source image datasets & Improved balanced accuracy and F1 score for REBA-based postural risk prediction. \\
\hline

44 & Xuke Yan \cite{yan2025efficient} (2025) & Efficient ROM analysis on edge devices & Reduced and quantised hourglass model for monocular ROM estimation & Custom ROM video data & Achieved 93.6\% joint prediction accuracy and RMSE of 3.21$^\circ$--4.25$^\circ$ on edge hardware. \\
\hline

45 & Brandon Schoenwether \cite{schoenwether2025reliability} (2025) & Reliability of markerless gait analysis & KinaTrax markerless motion capture with repeated-trial reliability analysis & 9 healthy adults over 3 sessions & Reported excellent reliability for most spatiotemporal parameters, particularly in the sagittal plane. \\
\hline

\end{longtable}}
\begingroup

\centering


{\scriptsize
\setlength{\tabcolsep}{4pt}
\renewcommand{\arraystretch}{1.08}
\begin{longtable}{c p{4.5cm} p{4.3cm} p{3.2cm}}
\caption{Comparison of classification accuracy reported by representative HAR-based studies on commonly used public datasets for gait, posture, mobility, and joint-movement analysis.}
\label{TableDataset_CA} \\
\hline
\textbf{SN} & \textbf{Study} & \textbf{Dataset} & \textbf{Classification Accuracy} \\
\hline

1 &
J. Wang \cite{wang2012mining} (2012) \newline
R. Vemulapalli \cite{vemulapalli2014human} (2014) \newline
Y. Du \cite{du2015hierarchical} (2015)
&
MSR-Action3D dataset
&
85.75\% \newline
92.46\% \newline
94.49\% \\
\hline

2 &
J. Wang \cite{wang2012mining} (2012) \newline
Q. Ke \cite{ke2017skeletonnet} (2017)
&
CMU MoCap dataset
&
98.13\% \newline
84.83\% \\
\hline

3 &
R. Vemulapalli \cite{vemulapalli2014human} (2014) \newline
M. A. R. Ahad \cite{ahad2021action} (2021)
&
UTKinect-Action dataset
&
97.08\% \newline
94.73\% \\
\hline

4 &
R. Vemulapalli \cite{vemulapalli2014human} (2014) \newline
M. A. R. Ahad \cite{ahad2021action} (2021)
&
Florence3D-Action dataset
&
90.88\% \newline
96.00\% \\
\hline

5 &
Y. Du \cite{du2015hierarchical} (2015)
&
Berkeley MHAD dataset
&
100.00\% \\
\hline

6 &
Y. Du \cite{du2015hierarchical} (2015)
&
HDM05 dataset
&
96.92\% \\
\hline

7 &
D. Leightley \cite{leightley2016automated} (2016)
&
K3Da dataset
&
87.10\% \\
\hline

8 &
Q. Ke \cite{ke2017skeletonnet} (2017)
&
SBU Kinect Interaction dataset
&
93.47\% \\
\hline

9 &
Q. Ke \cite{ke2017skeletonnet} (2017) \newline
C. Caetano \cite{caetano2019skeleton} (2019) \newline
T. Ahmad \cite{ahmad2019action} (2019)
&
NTU RGB+D dataset
&
81.16\% \newline
80.30\% \newline
90.70\% \\
\hline

10 &
T. Ahmad \cite{ahmad2019action} (2019)
&
Stanford 40 Actions dataset
&
84.80\% \\
\hline

11 &
T. Ahmad \cite{ahmad2019action} (2019)
&
PKU-MMD dataset
&
95.50\% \\
\hline

12 &
Y. Li \cite{li2020human} (2020) \newline
X. Lv \cite{lv2022analysis} (2022)
&
COCO2017 dataset
&
88.30\% \newline
97.19\% \\
\hline

13 &
C. Fan \cite{fan2020gaitpart} (2020)
&
CASIA-B dataset
&
88.90\% \\
\hline

14 &
C. Fan \cite{fan2020gaitpart} (2020)
&
OU-MVLP dataset
&
88.70\% \\
\hline

15 &
M. A. R. Ahad \cite{ahad2021action} (2021)
&
MSR 3D Action Pairs dataset
&
95.45\% \\
\hline

16 &
M. A. R. Ahad \cite{ahad2021action} (2021)
&
Online Action Detection (OAD) dataset
&
98.44\% \\
\hline

17 &
M. A. R. Ahad \cite{ahad2021action} (2021) \newline
A. Suhendar \cite{suhendar2022skeletal} (2022)
&
KARD dataset
&
98.11\% \newline
99.40\% \\
\hline

\end{longtable}
}

\endgroup

\section*{Abbreviations}\addcontentsline{toc}{section}{Abbreviations}

{\scriptsize
\setlength{\tabcolsep}{6pt}
\renewcommand{\arraystretch}{1.12}
\setlength\LTleft{\fill}
\setlength\LTright{\fill}

\begin{longtable}{p{3cm} p{10.5cm}}
\caption{List of abbreviations used in the paper.}
\label{tab:abbreviations}\\
\hline
\textbf{Term} & \textbf{Represent} \\
\hline
\endfirsthead

\hline
\textbf{Term} & \textbf{Represent} \\
\hline
\endhead

JMP & Joint Mobility Problem / Joint Mobility Problems \\
HAR & Human Action Recognition \\
HPE & Human Pose Estimation \\
CV & Computer Vision \\
GA & Gait Analysis \\
OA & Osteoarthritis \\
ROM & Range of Motion \\
KAM & Knee Adduction Moment \\
TUG & Timed Up and Go \\
REBA & Rapid Entire Body Assessment \\
BMI & Body Mass Index \\
IMU & Inertial Measurement Unit \\
MoCap & Motion Capture \\
RGB & Red, Green, Blue \\
RGB-D & Red, Green, Blue and Depth \\
CNN & Convolutional Neural Network \\
RNN & Recurrent Neural Network \\
LSTM & Long Short-Term Memory \\
ConvLSTM & Convolutional Long Short-Term Memory \\
ConvRNN & Convolutional Recurrent Neural Network \\
CNNRNN & Convolutional Neural Network--Recurrent Neural Network \\
GCN & Graph Convolutional Network \\
ANN & Artificial Neural Network \\
SVM & Support Vector Machine \\
DTW & Dynamic Time Warping \\
HMM & Hidden Markov Model \\
LDA & Linear Discriminant Analysis \\
PCA & Principal Component Analysis \\
ROC & Receiver Operating Characteristic \\
SPM & Statistical Parametric Mapping \\
GMM & Gaussian Mixture Model \\
BPTT & Back-Propagation Through Time \\
MAE & Mean Absolute Error \\
RMSE & Root Mean Square Error \\
MSE & Mean Squared Error \\
RMSD & Root Mean Square Difference \\
ICC & Intraclass Correlation Coefficient \\
PCK & Percentage of Correct Keypoints \\
AP & Average Precision \\
FPS & Frames Per Second \\
HMI & Human--Machine Interaction \\
PAF & Part Affinity Fields \\
EEMD & Ensemble Empirical Mode Decomposition \\
IMF & Intrinsic Mode Function \\
IFC & Intensive Feature Consistency \\
RQ-HG & Reduced and Quantized Hourglass \\
HRNet & High-Resolution Network \\
HRNN & Hierarchical Recurrent Neural Network \\
HBRNN-L & Hierarchical Bidirectional Recurrent Neural Network with LSTM \\
HURNN-L & Hierarchical Unidirectional Recurrent Neural Network with LSTM \\
DBRNN-L & Deep Bidirectional Recurrent Neural Network with LSTM \\
DURNN-L & Deep Unidirectional Recurrent Neural Network with LSTM \\
DBRNN-T & Deep Bidirectional Recurrent Neural Network with \texttt{tanh} units \\
DURNN-T & Deep Unidirectional Recurrent Neural Network with \texttt{tanh} units \\
TSRJI & Tree Structure Reference Joints Image \\
BMLD & Bi-gram with Maximum Likelihood Decoding \\
NERF & Non-Euclidean Relational Fuzzy \\
FTP & Fourier Temporal Pyramid \\
LOP & Local Occupancy Pattern \\
LJPF & Linear Joint Position Feature \\
AJPF & Angular Joint Position Feature \\
KPF & Kinematic Posture Feature \\
CO & Chain Order \\
BODY\_25 / B25 & OpenPose 25-body-joint pose model \\
COCO & Common Objects in Context \\
COCO2017 & Common Objects in Context 2017 dataset \\
KARD & Kinect Activity Recognition Dataset \\
OAD & Online Action Detection dataset \\
UTKinect & University of Tennessee Kinect Action dataset \\
MHAD & Multimodal Human Action Database \\
HDM05 & Human Dynamics Database 2005 \\
PKU-MMD & Peking University Multi-Modality Dataset \\
CASIA-B & Chinese Academy of Sciences Institute of Automation Dataset B \\
OU-MVLP & Osaka University Multi-View Large Population dataset \\
CMU MoCap & Carnegie Mellon University Motion Capture dataset \\
NTU RGB+D & Nanyang Technological University RGB+D dataset \\
MSR-Action3D & Microsoft Research Action3D dataset \\
MSRDailyActivity3D & Microsoft Research Daily Activity 3D dataset \\
Human3.6M & Human3.6M pose dataset \\
ENSAM & \mbox{\'E}cole Nationale Sup\'erieure d'Arts et M\'etiers gait dataset \\
VICON & Vicon motion capture system \\
KneeKG & Knee Kinesiography system \\
KSS & Knee Society Score \\
WOMAC & Western Ontario and McMaster Universities Osteoarthritis Index \\
HOOS & Hip disability and Osteoarthritis Outcome Score \\
HHS & Harris Hip Score \\
SF-36 & 36-Item Short Form Health Survey \\
KL & Kellgren--Lawrence grading scale \\
THR & Total Hip Replacement \\
TKA & Total Knee Arthroplasty \\
TKR & Total Knee Replacement \\
RTHA & Resurfacing Total Hip Arthroplasty \\
STHA & Standard Total Hip Arthroplasty \\
GJH & Generalized Joint Hypermobility \\
ACL & Anterior Cruciate Ligament \\
ML & Machine Learning \\
\hline
\end{longtable}
}

\begingroup\small
\bibliographystyle{unsrtnat}
\bibliography{bibliography}
\endgroup
\end{document}